% Modernised reading — fonds Alexandre Grothendieck, shelfmark 156-8, pages 1-126.
%
% This is an INTERPRETATION, not a transcription. Everything the critical
% apparatus carried moves into footnotes. In any dispute of substance the
% transcription governs: whoever cites Grothendieck must cite batch-NN.fr.tex.
%
% Pass: Opus 5.5 (claude-opus-5-5), 2026-10-10 — first pass, unchecked against
% the pages by a human. The folder was read whole: seven batches, pages 1-126.
% Pages 4 and 6 are typed versos of another text and were not transcribed;
% everything else is his mathematics. Dated 26 June 1986 (page 7), 27 June
% (page 9), « 27 juin (ou plutôt, 28 juin !) » (page 18), 1 July (page 89),
% 2 July (page 97), 4 July (pages 102 and 117); a margin of page 36 reads
% 11.7.86 (or 1.7.86). Not measured against the Opus 5 reference reading of
% folder 115.
%
% What the folder is. Chapter VIII (« GF VIII ») of Vers une géométrie des
% formes, the fourth « mouture » of the Analysis situs (after GF IV = 156-4,
% GF VI = 156-6, GF VII = 156-7), followed by GF IX (156-9, from 5 July). It
% restarts the axiomatics from the « magasin » that the last pages of GF VII
% (156-7 pp. 109-110) had just reached: a set M of « multistructures » with
% three relations — incidence ≤, raffinement ≪ (then the interior incidence
% ≪̊ as primitive, pp. 28-31), interior disjointness |∘| — and builds on it
% figures, ateliers, lieux, support functions, interior supports (a polarity),
% and finally works out everything for the magasin of an ordered set (points
% and intervals, pp. 89-126). Pages 1-6 are a cover, a later table of the
% hierarchy of magasins, and an unfinished aside on oriented magasins.
%
% Pagination. His own runs 1, 2, 1bis, 2bis, 3, ..., 118 on pages 7-126:
% his page n is our page n+8 for n ≥ 3 (1, 2, 1bis, 2bis = 7, 8, 9, 10). His
% internal references are converted in the text. The pencilled archive numbers
% on the leaves run one ahead from page 7 on; we follow the facsimile as served.
%
% Notation fixed for the whole document. M the magasin, X, Y, Z its elements;
% ≤ incidence, X̃ = M_{≤X}, ∂X = X̃ ∖ {X}; ≪ raffinement; ≪̊ interior
% incidence; |∘| interior disjointness; ≶ compatibility (pages 81-96 of the
% transcription set it as a barred ≶, read there \not\lessgtr: same sign); ∥
% disjunction. Fig(M) the set of figures (𝓕_M on page 8), 𝔉 an atelier
% (« atlas » on pages 61-67: same notion). ℒ the lieux; Omb(X) = M_{≪X},
% Omb°(X) = M_{≪̊X}, omb = Omb ∩ ℒ, omb° = Omb° ∩ ℒ. cosupp°, supp°, Σ_M.
% A support function X ↦ |X| ⊂ P, |X|° = |X| ∖ |∂X|, φ(A) = ∪_{X∈A}|X|°.
% For an ordered set L: P_x, S_{a,b} (a < b), M_0, M_1, δ, Drap(L) (finite
% chains), Drap*(L), F_T, <_pos, <¹. « Préfigure » throughout (pages 102-112
% write mostly « profigure »). TWO BOUNDARIES: ∂_c Φ (union of the boundaries
% of the connected components, pages 103-105) and ∂Φ (vertices of order 1,
% margin of page 104 and pages 118-126); they differ exactly at lacunary
% vertices, and the reading keeps them apart.
%
% Substantive departures from the pages, each footnoted where it occurs:
% — p. 9 (margin): « X ≠ Y ⇒ (X |∘| Y ⇔ X ≶ Y) »: only ⇐ holds; counter-example.
% — p. 10, Cor. 3: read with ≪ (X ≶ Y ⇒ (X ≪ Y ⇔ X ≤ Y)); with ≪̊ as the
%   transcription has it the statement fails for X < Y.
% — p. 11, Prop. 3: G ∩ Omb(F') is the greatest sub-figure of G refining F';
%   commutation with Inf uses G ≪ F; proof supplied.
% — p. 12, Prop. 4 (gluing of refinements) and p. 12 Prop. 5: stated without
%   proof on the page; not proved here.
% — p. 23: « prépolyidéal » read as « two majorised elements have no common
%   lower bound or an Inf », the only reading under which p. 24's equivalence
%   holds.
% — p. 30: Mag 4' carries a doubtful « non »; dropped (it contradicts Mag 4);
%   the form restricted to comparable elements is weaker than Mag 4.
% — p. 30: « quinze semaines » (uncertain) — the dates give about a fortnight.
% — p. 62 and p. 76: « strict. dirigé » (uncertain) read « strictement
%   divisible » (his p. 23 = our p. 31, Mag div).
% — p. 43: « ≠ » read « = ∅ »; p. 50: « vide » read « non vide »; p. 53:
%   A_{i-1} read A_{i+1}; p. 60: ≪_M | M read ≪_M | M'; p. 83, p. 85, p. 89,
%   p. 91 misprints corrected as noted.
% — pp. 45-50: the proposition on indexed families needs i ≤ j ⇒ X_i ⊂ X_j
%   (his margin p. 45 and his own query p. 47); made a standing hypothesis.
% — pp. 76-77: the bijection Σ_ℒ ≅ Σ_M (« il faudra bien que je me tape la
%   vérification ! ») is verified here; p. 76's Omb(A) read Omb°(A).
% — p. 88: « Omb is a fully faithful support function in a faithful
%   magasin » is false (the magasin of intervals of ℝ); the axiom of supports
%   does follow from the existence of SOME fully faithful support function.
% — p. 102: his computation supp°(P_a) = {P_a, P_b} (« Canular ») is wrong;
%   supp°(P_a) = {P_a}; what stands is supp°(S̃_{a,b}) = M ≠ Omb(S_{a,b}), in
%   an L violating Hyp. H.
% — p. 112: « supp Φ » read supp° Φ.
% — p. 126: Cosupp°(Φ) = Omb°(Φ'*) ∪ M_{<a} ∪ M_{>b}, with the whole Φ'*
%   (points included); with Φ'*_1 alone (as the page has it, the index
%   rewritten) the formula misses Δ ∖ Φ_0.
%
% Modern names given and footnoted as ours: graphe au sens de Serre (1977),
% propriété du diamant, idéal d'un ensemble ordonné, éléments complètement
% sup-irréductibles, complexe polyédral / complexe CW régulier, méréotopologie
% (Whitehead 1929, Clarke 1981), polarité de Birkhoff (1940), espace
% d'orthogonalité et orthotreillis complet (Dacey 1968), formule de
% projection, ordre d'intervalles (Fishburn 1970), complexe des chaînes,
% ensembles constructibles.
%
% What the folder announces and never establishes: Prop. 4 (gluing of
% refinements) and Prop. 5 (pp. 12-13); the programme of p. 13 (subdivisions);
% the comparison of the magasin and atelier points of view beyond the Scholie
% of p. 20; condition d) of morphisms (p. 36, queried); the hope of p. 78 that
% Σ_M sits in 𝔓(P) stable under Sup, Inf and complement (true only for
% admissible families, pp. 85-87); points 2° and 3° of p. 88 (magasins given by
% ≤ and |∘| alone; a failure of the axiom of supports — margin « Oui ! »); the
% closure statement 2) of p. 117 (« Vérifier ! »); the formula of p. 126,
% where the folder stops (it is taken up in GF IX, 156-9).
\documentclass[11pt,a4paper]{article}
\input{../preamble/grothendieck}

\folder{156-8}
\pages{1}{126}
\dating{1986}
\watermark{Édition de démonstration\\interprétation personnelle de l'œuvre}
\foldertitle{[Chapitre] VIII. Analysis situs (quatrième mouture) : notes manuscrites (26/06-04/07/1986) — lecture modernisée du dossier entier}

\begin{document}

\section*{Résumé}

\begin{resume}
Comment décrire une forme --- un polyèdre, une surface découpée en morceaux,
un dessin fait de points et de traits --- sans parler de ses points ? Ces cent
vingt-six pages, écrites du 26 juin au 4 juillet 1986, essaient de le faire en ne
retenant que les \emph{morceaux} et trois manières dont deux morceaux peuvent
se tenir l'un par rapport à l'autre : l'un est une face de l'autre (une arête
est une face d'un triangle) ; l'un est logé à l'intérieur de l'autre (un point
au milieu d'un segment) ; leurs intérieurs ne se rencontrent pas (deux
triangles qui ont un côté commun). Grothendieck appelle \emph{magasin}
l'ensemble des morceaux muni de ces trois relations, et \emph{atelier}
l'ensemble des assemblages qu'on en tire, les \emph{figures} : des morceaux
deux à deux compatibles, pris avec toutes leurs faces.

C'est la quatrième « mouture » de ce chapitre d'\emph{Analysis situs} en trois
semaines, et la page 30 en fait elle-même l'histoire : on est parti des
figures et de leurs relations, on est revenu aux morceaux, on a remplacé la
compatibilité par la disjonction des intérieurs --- « l'innovation
d'avant-hier » ---, et l'on prend enfin pour donnée primitive l'incidence
intérieure : « aujourd'hui, l'épanouissement intégral ! », et plus haut :
« Décidément, c'est la forme parfaite de définition ! ». Le lecteur assiste à
une pensée qui se corrige d'un jour sur l'autre, et qui dit à chaque fois
pourquoi le nouveau choix est meilleur : plus économe pour écrire les axiomes,
plus proche de l'intuition, mieux séparé dans ses fonctions.

L'idée directrice vient de la façon dont on connaît une région sur une carte
quand on n'a pas le droit de nommer ses points : par ce qui ne la chevauche
pas. Ce qui ne chevauche rien de ce qui ne chevauche pas $A$ est le
\emph{support} de $A$ --- une double négation, comme l'orthogonal de
l'orthogonal en algèbre linéaire. Les plus petits morceaux, les \emph{lieux},
jouent le rôle des points ; une \emph{fonction support} réalise les morceaux
comme de vraies parties d'un ensemble, et la question est de savoir quand la
double négation abstraite et la géométrie ensembliste disent la même chose.
Réponse : toujours pour les familles de morceaux tirées d'une même figure,
où « toute pathologie disparaît ».

Le dernier tiers est un banc d'essai, le plus simple possible : la droite,
ou plus généralement un ensemble ordonné, découpé en points et en intervalles.
Tout y est calculé --- les figures sont des chaînes de chaînes, les supports
des réunions d'intervalles ouverts, fermés ou semi-ouverts --- jusqu'à une
classification des sommets d'une « préfigure » en cinq espèces, isolés,
de bord propres ou impropres, redondants, lacunaires. Le dossier s'arrête au
milieu d'une formule pour le complémentaire d'une préfigure ; le chapitre IX
la reprend le lendemain.

On reconnaîtra, sous d'autres noms, l'ensemble ordonné des faces d'un complexe
cellulaire, la théorie des régions qu'on appelle méréotopologie, les
correspondances de Galois et les orthotreillis, les ordres d'intervalles. Rien
n'indique que Grothendieck ait connu ces travaux ; il construit, ligne à
ligne, ce dont il a besoin, et s'arrête quand il ne sait pas : « Il faudra
bien que je me tape la vérification ! ».

\keywords{face poset, regular CW complex, polyhedral complex, diamond property, Serre graph, order ideal, completely join-irreducible element, mereotopology, Galois connection, orthogonality space, complete ortholattice, projection formula, interval order, order complex, constructible set}
\end{resume}

\section*{Le fil du dossier, et les conventions}

\begin{enumerate}
\item[1.] \emph{Couverture, table des magasins, magasins orientés} (pages 1 à
6) : une table récapitulative écrite après coup, et un essai inachevé.
\item[2.] \emph{Les données d'un magasin} (pages 7 et 8), le 26 juin :
axiomes Mag 1 à Mag 4, compatibilité, disjonction, figures.
\item[3.] \emph{Premières conséquences} (pages 9 et 10), le 27 juin.
\item[4.] \emph{Restriction et recollement des raffinements} (pages 11 à 13).
\item[5.] \emph{Lieux, et les magasins divisibles, quasi-ensemblistes,
ensemblistes} (pages 14 à 17).
\item[6.] \emph{Ateliers} (pages 18 à 23), le 28 juin.
\item[7.] \emph{La variante polyidéale} (pages 23 à 28).
\item[8.] \emph{L'incidence intérieure comme donnée primitive} (pages 28 à 32).
\item[9.] \emph{Magasins triviaux et figures indexées} (pages 33 à 35).
\item[10.] \emph{Morphismes de magasins} (pages 35 à 40).
\item[11.] \emph{Fonction support} (pages 41 à 44).
\item[12.] \emph{Familles indexées de parties} (pages 45 à 53).
\item[13.] \emph{Images inverses, quotients, sous-magasins} (pages 53 à 60).
\item[14.] \emph{Homomorphismes d'ateliers} (pages 61 à 66).
\item[15.] \emph{Fonctions support exactes} (pages 66 à 72).
\item[16.] \emph{Supports intérieurs} (pages 73 à 80).
\item[17.] \emph{Supports constructibles et pleine fidélité} (pages 81 à 88).
\item[18.] \emph{Le magasin d'un ensemble ordonné} (pages 89 à 96), le
premier juillet.
\item[19.] \emph{Figures et supports élémentaires} (pages 97 à 102), le 2
juillet.
\item[20.] \emph{Préfigures} (pages 102 à 110), le 4 juillet.
\item[21.] \emph{Supports des préfigures} (pages 110 à 117).
\item[22.] \emph{Sommets, et la description des préfigures} (pages 117 à 126).
\end{enumerate}

\emph{La pagination.} À partir de la page 7, les feuillets portent sa propre
pagination 1, 2, 1bis, 2bis, 3, \ldots, 118. Sa page $n$ est notre page
$n + 8$ pour $n \geq 3$ ; ses pages 1, 2, 1bis, 2bis sont nos pages 7, 8, 9,
10. Ses renvois internes sont convertis ci-dessous, avec son numéro entre
parenthèses quand il importe. Les pages 4 et 6 sont des versos dactylographiés
d'un autre texte, non transcrits.

\emph{La place du dossier.} C'est le chapitre VIII de \emph{Vers une géométrie
des formes}, la quatrième mouture de l'\emph{Analysis situs} après les
chapitres IV (dossier 156-4), VI (156-6) et VII (156-7). Le mot « magasin » et
la relation $\mathrel{|\circ|}$ de disjonction intérieure sont apparus aux
dernières pages du chapitre VII (pages 101 à 110 du dossier 156-7) ; les
figures ensemblistes viennent du chapitre V, « Algèbre des figures » (156-5).
Le chapitre IX (156-9), commencé le 5 juillet, reprend le magasin d'un
ensemble ordonné là où celui-ci s'arrête.

\emph{Conventions, valables pour tout le dossier.} Un \emph{magasin} est un
ensemble $\mathcal{M}$ dont les éléments $X, Y, Z$ sont appelés
\emph{multistructures}, muni de :
\begin{itemize}
\item[--] l'\emph{incidence} $X \leq Y$ (« $X$ est une strate de $Y$ ») ; on
note $\widetilde{X} = \mathcal{M}_{\leq X} = \{Y \mid Y \leq X\}$ et
$\partial X = \widetilde{X} \smallsetminus \{X\}$ ;
\item[--] le \emph{raffinement} $X \ll Y$ (« $X$ est contenu dans $Y$ »), et
l'\emph{incidence intérieure} $X \mathrel{\mathring{\ll}} Y$ (« $X$ est dans
l'intérieur de $Y$ ») ;
\item[--] la \emph{disjonction intérieure} $X \mathrel{|\circ|} Y$ (« les
intérieurs de $X$ et $Y$ sont disjoints »).
\end{itemize}
On en déduit la \emph{compatibilité} $X \lessgtr Y$ et la \emph{disjonction}
$X \parallel Y$ (page 7)\footnote{Le signe de compatibilité est, sur la page,
un $<$ et un $>$ croisés, parfois posés sur un trait ; la transcription le rend
$\lessgtr$, et sur les pages 81 à 96 $\not\lessgtr$. C'est partout la même
relation, notée ici $\lessgtr$.}. $\mathrm{Fig}(\mathcal{M})$ est l'ensemble des
figures (noté $\mathcal{F}_{\mathcal{M}}$ page 8), $\mathfrak{F}$ un atelier.
$\mathcal{L}$ est l'ensemble des lieux ; $\mathrm{Omb}(X) =
\mathcal{M}_{\ll X}$ et $\mathrm{Omb}^{\circ}(X) = \mathcal{M}_{\mathring{\ll}
X}$ sont l'ombre et l'ombre intérieure, $\mathrm{omb}$ et $\mathrm{omb}^{\circ}$
leurs traces sur $\mathcal{L}$. Une fonction support est notée $X \mapsto |X|
\subset P$.

\emph{Deux homonymies, fixées ici.} Les pages 61 à 67 appellent « atlas » ce
que le reste du dossier appelle « atelier » ; c'est la même notion, et l'on dit
partout atelier. Plus sérieusement, le \emph{bord} d'une préfigure reçoit deux
définitions qui ne coïncident pas : aux pages 103 à 105 c'est la réunion des
bords de ses composantes connexes, que l'on note ici $\partial_{c}\Phi$ ; en
marge de la page 104 puis aux pages 118 à 126, c'est l'ensemble des sommets
d'où part exactement un intervalle de $\Phi$, noté $\partial\Phi$. Elles
diffèrent exactement aux sommets « lacunaires » (page 119), et la suite du
dossier est écrite avec la seconde.

\emph{Ce que le dossier annonce et n'établit pas} : le recollement des
raffinements et la proposition 5 des pages 12 et 13 ; le programme de la page
13 sur les subdivisions ; la condition d) des morphismes, mise en question par
lui-même (page 36) ; l'espoir de la page 78 de plonger tous les supports dans
les parties d'un ensemble en respectant Sup, Inf et complémentaire (vrai
seulement pour les familles admissibles, pages 85 à 87) ; les points 2° et 3°
du programme de la page 88 ; l'énoncé de fermeture de la page 117
(« Vérifier ! ») ; la formule de la page 126, où le dossier s'arrête.

\pagerange{1}{6}
\section*{Couverture, table des magasins, magasins orientés (pages 1 à 6)}

La couverture (page 1) porte, de sa main, « GF VIII », et renvoie à une page
de numéro douteux pour la « définition des magasins »\footnote{« Voir page 46
Déf de magasins », le nombre lu avec doute. Ni sa page 46 (notre page 54) ni
notre page 46 ne contiennent cette définition ; elle est aux pages 7 et 30.}.

\subsection*{La table de la page 2}

La page 2 dresse une hiérarchie des sortes de magasins :

\begin{tikzcd}[column sep=tiny, row sep=small, nodes={font=\scriptsize}]
 & \text{généraux} \arrow[dl, no head] \arrow[dr, no head] & \\
\text{locaux} \arrow[dr, no head] & & \text{fidèles} \arrow[dl, no head] \arrow[d, no head] \\
 & \text{ponctuaires} \arrow[d, no head] & \text{maquettes} \\
 & \text{ensemblistes} \arrow[d, no head] & \\
 & \text{modérés} &
\end{tikzcd}

Elle a été écrite après coup : elle renvoie à ses pages 9 à 31 (nos pages 17 à
39) et utilise l'ensemble $\Sigma_{\mathcal{M}}$ des supports, qui n'apparaît
qu'à la page 75. Ses noms ne sont pas ceux du texte, et la correspondance est
la suivante\footnote{Établie d'après les numéros de page que donne la table.
« Magasins locaux (p.~23) » renvoie à notre page 31, les magasins strictement
divisibles ; « ponctuaires (p.~24) » à notre page 32, les quasi-ensemblistes ;
« maquettes (p.~25) » à notre page 33, les magasins triviaux ; « fidèles
(p.~31) » à notre page 39 ; « ensemblistes (p.~9) » à notre page 17.} :
\begin{itemize}
\item[--] \emph{généraux} : $(\mathcal{M}, \leq, \mathring{\ll},
\mathrel{|\circ|})$ soumis à Mag 1 à Mag 4 (pages 7 et 30) ;
\item[--] \emph{locaux} : les magasins strictement divisibles (pages 15 et 31),
où $\mathrel{|\circ|}$ se lit sur les lieux ; la note en marge dit qu'un tel
magasin « se réfère à des voisinages particuliers dévolus aux lieux », et que
la trace sur les lieux ne respecte ni $\cap$ ni $\bigcup$ ;
\item[--] \emph{ponctuaires} : les quasi-ensemblistes (pages 16 et 32), pour
lesquels $\Sigma_{\mathcal{M}} \simeq \mathfrak{P}(\mathcal{L})$ (page 77) ;
un lieu $x$ y est \emph{ponctuel} si $x \mathrel{|\circ|} X \Leftrightarrow x
\notin \mathrm{Omb}^{\circ}(X)$ pour tout $X$, et ce sont, dit la marge, « les
seuls magasins qui correspondent plus ou moins » à ce qu'on admet communément
pour la topologie ;
\item[--] \emph{fidèles} : $(\mathcal{M}, \leq, \mathring{\ll})$, la
disjonction intérieure étant définie par $\mathring{\ll}$ (page 39), de sorte
que $X \mapsto \mathrm{Omb}^{\circ}(X)$ est une fonction support fidèle ;
\item[--] \emph{maquettes} : un ensemble ordonné $(\mathcal{M}, \leq)$ seul,
$\mathring{\ll}$ étant l'égalité (page 33) ; la maquette d'un espace $T_{0}$
n'est ponctuaire que s'il est discret ;
\item[--] \emph{ensemblistes} : $\mathcal{M} \subset \mathrm{Figél}(L)$ soumis
à Magens 1 et 2 (page 17), et \emph{modérés} avec un troisième axiome, qui
« correspondent à des topologies modérées, qui n'ont été développées à présent
que pour la seule structure linéaire par morceaux » ; la marge ajoute qu'il y a
des magasins modérés non ensemblistes, les maquettes\footnote{La « topologie
modérée » est celle que réclamait l'\emph{Esquisse d'un programme} (1984) ;
voir la lecture du dossier 156-2, page 2.}.
\end{itemize}
Une dernière note change les notations : $\trianglelefteq$ pour l'incidence,
$\lhd$ pour l'incidence stricte, $\subseteq$ pour l'inclusion des figures.
Ces notations ne sont pas reprises dans le corps du dossier, qui garde $\leq$.

\subsection*{Un essai de magasins orientés (pages 3 et 5)}

Les pages 3 et 5 esquissent des « multistructures orientées ». Les données sont
une incidence immédiate $\lhd$, l'incidence intérieure, la disjonction
intérieure, une involution sans point fixe $X \mapsto X^{-}$ (« la
multistructure orientée opposée ») et une partie $\mathcal{M}_{0}^{+}$ de
$\mathcal{M}$. Les axiomes Magor demandent que $\lhd$ engendre un ordre strict
sans chaîne fermée satisfaisant la condition des chaînes descendantes (tout
élément est de dimension finie), que l'involution respecte les relations, que
$X < Y$ entraîne $X \mathrel{|\circ|} Y$, que les éléments de
$\mathcal{M}_{0}^{+}$ soient de codimension nulle et que les éléments minimaux
se répartissent en $\mathcal{M}_{0} = \mathcal{M}_{0}^{+} \amalg
(\mathcal{M}_{0}^{+})^{-}$. Deux conditions sont proprement orientées :
\begin{enumerate}
\item[a)] si $Z < X$ est de codimension 2, il y a un seul $Y^{\sharp}$ avec $Z
\lhd Y^{\sharp} \lhd X$, et un $Y^{\flat}$ du côté des orientations
opposées\footnote{La lecture des signes de cette condition est incertaine, et
une flèche de la page la déplace avant Magor 4 sans qu'on voie où. Elle
rappelle la propriété dite aujourd'hui « du diamant » de l'ensemble des faces
d'un polytope ou d'un complexe régulier --- entre deux faces de codimension 2
il y en a exactement deux ---, sous la forme orientée qui donne $\partial \circ
\partial = 0$ en homologie cellulaire. Le rapprochement est le nôtre ; nous ne
reconstruisons pas les signes.} ;
\item[b)] une multistructure de dimension 1 a exactement un sommet dans
$\mathcal{M}_{0}^{+}$.
\end{enumerate}
Lorsque $\mathring{\ll}$ est l'égalité et $\mathrel{|\circ|}$ la différence
(« maquettes orientées ») et que tout est de dimension $\leq 1$, on trouve un
ensemble $S = \mathcal{M}_{0}^{+}$ de sommets, un ensemble $A$ d'arêtes
orientées muni d'une involution sans point fixe, et une application $A \to S$,
l'origine : ce sont, dit la page, « les graphes ». C'est exactement la
définition des graphes de Serre\footnote{J.-P. Serre, \emph{Arbres, amalgames,
$\mathrm{SL}_{2}$} (1977), chapitre I, § 2 : un graphe est un ensemble de
sommets, un ensemble d'arêtes orientées muni d'une involution sans point fixe
$e \mapsto \bar{e}$, et une application origine. L'extrémité d'une arête est
l'origine de l'arête opposée, ce que donne ici la condition b) appliquée à
$X^{-}$. La page ne cite personne.}. En dimension 2, la page annonce « la
description des complexes cellulaires combinatoires de dimension $\leq 2$ », et
s'arrête : la page 7 reprend tout à neuf.

\pagerange{7}{8}
\section*{Les données d'un magasin (pages 7 et 8)}

Le 26 juin, sous le titre « Analysis Situs (quatrième mouture) », la page 7
« reprend le formalisme algébrique » du chapitre VII, page 110\footnote{Dossier
156-7, page 110, où les axiomes Mag 1 à Mag 4 sont écrits au propre pour la
première fois. La note de la page 7 dit que Mag 4 « assouplit légèrement »
l'ancien At 2.}.

\begin{quote}
\textbf{Définition} (page 7). Un \emph{magasin} est un ensemble $\mathcal{M}$
muni de trois relations $\leq$, $\ll$, $\mathrel{|\circ|}$ telles que :
\end{quote}
\begin{enumerate}
\item[Mag 1.] $\leq$ et $\ll$ sont des relations d'ordre, et $X \leq Y
\Rightarrow X \ll Y$.
\item[Mag 2.] Si $X \ll Y$, l'ensemble $\widetilde{Y}^{X} = \{Z \in
\widetilde{Y} \mid X \ll Z\}$ a un plus petit élément. On écrit $X
\mathrel{\mathring{\ll}} Y$ quand ce plus petit élément est $Y$ lui-même.
\item[Mag 3.] $\mathrel{|\circ|}$ est symétrique et antiréflexive, et
$X \mathrel{|\circ|} Y$, $X' \mathrel{\mathring{\ll}} X$, $Y'
\mathrel{\mathring{\ll}} Y$ entraînent $X' \mathrel{|\circ|} Y'$.
\item[Mag 4.] Si $Y, Z \leq X$ et $Y \neq Z$, alors $Y \mathrel{|\circ|} Z$.
\end{enumerate}


Le modèle à garder en tête est l'ensemble des cellules fermées d'un complexe
cellulaire : $X \leq Y$ si $X$ est une face de $Y$, $X \ll Y$ si $X$ est
contenue dans $Y$ comme partie, $X \mathrel{\mathring{\ll}} Y$ si l'intérieur
de $X$ est dans l'intérieur de $Y$, $X \mathrel{|\circ|} Y$ si leurs intérieurs
sont disjoints. Mag 2 dit qu'une cellule contenue dans $Y$ a son intérieur dans
une seule face de $Y$, la plus petite qui la contienne ; Mag 4, que deux faces
distinctes d'une même cellule ont des intérieurs disjoints.

Deux remarques que la page laisse implicites, et dont tout ce qui suit se sert.
La relation $\mathring{\ll}$ est réflexive (si $Z \leq X$ et $X \ll Z$, alors
$Z \ll X \ll Z$, donc $Z = X$), et elle est contenue dans $\ll$. Et $X \ll Y$
équivaut à l'existence d'un $Y' \leq Y$ avec $X \mathrel{\mathring{\ll}} Y'$
--- le plus petit élément de Mag 2 ---, ce que la page 28 prendra pour
définition de $\ll$.

\begin{quote}
\textbf{Définitions} (page 7). On pose
\[
\begin{aligned}
X \lessgtr Y &\iff X' \mathrel{|\circ|} Y' \text{ pour tous } X' \in
\widetilde{X} \smallsetminus \widetilde{Y},\ Y' \in \widetilde{Y}
\smallsetminus \widetilde{X} ; \\
X \parallel Y &\iff X' \mathrel{|\circ|} Y' \text{ pour tous } X' \in
\widetilde{X},\ Y' \in \widetilde{Y}.
\end{aligned}
\]
\end{quote}

Deux multistructures sont \emph{compatibles} quand elles peuvent coexister
dans une même figure, se touchant au plus le long de faces communes ; elles
sont \emph{disjointes} quand elles ne se touchent pas du tout. On a $X \parallel
Y$ si et seulement si $X \lessgtr Y$ et $\widetilde{X} \cap \widetilde{Y} =
\emptyset$ : une strate commune serait disjointe d'elle-même.

\begin{quote}
\textbf{Définition} (page 8). Une \emph{figure} est une partie $F \subset
\mathcal{M}$ fermée pour $\leq$ (si $X \in F$ et $Y \leq X$, alors $Y \in F$)
et formée d'éléments deux à deux compatibles.
\end{quote}

Ce sont les sous-complexes, dans le modèle cellulaire. La page envisage deux
restrictions éventuelles : des conditions de finitude, et une condition
« polyédrale » --- si $X, Y \in F$ ont une strate commune, $\widetilde{X} \cap
\widetilde{Y}$ a un plus grand élément --- qui sera la variante polyidéale des
pages 23 à 28. Elle appelle l'ensemble $\mathrm{Fig}(\mathcal{M})$ « l'édition
maximale » associée au magasin\footnote{Lecture incertaine de la marge.}.

$\mathcal{M}$ se plonge dans $\mathrm{Fig}(\mathcal{M})$ par $X \mapsto
\widetilde{X}$, la \emph{figure élémentaire} de $X$, et les relations
s'étendent aux figures :
\[
\begin{aligned}
F \leq G &\iff F \subset G, \\
F \ll G &\iff \text{tout } X \in F \text{ raffine un } Y \in G, \\
F \lessgtr G &\iff X \lessgtr Y \text{ pour tous } X \in F,\ Y \in G
\iff F \cup G \in \mathrm{Fig}(\mathcal{M}), \\
F \parallel G &\iff F \lessgtr G \text{ et } F \cap G = \emptyset .
\end{aligned}
\]
Les bornes inférieures pour $\leq$ sont les intersections (de familles non
vides) ; la borne supérieure de $F$ et $G$ existe si et seulement si $\{F,G\}$
est majorée, si et seulement si $F \lessgtr G$, et c'est alors $F \cup G$. On
n'écrit $\mathrel{|\circ|}$ et $\mathring{\ll}$ qu'entre éléments de
$\mathcal{M}$ (page 11).

\pagerange{9}{10}
\section*{Premières conséquences (pages 9 et 10)}

Le 27 juin, deux pages numérotées 1bis et 2bis, que la marge de la page 8
demande d'insérer avant les sorites sur les figures, donnent « quelques
compléments sur $\mathcal{M}$ lui-même ». Le bas de la page 9 est un
palimpseste de tentatives biffées ; voici ce qui en résulte.

\begin{quote}
\textbf{Proposition $1_{0}$.} a) $\lessgtr$ est symétrique et
réflexive. b) Si $X' \leq X$, $Y' \leq Y$ et $X \lessgtr Y$, alors $X' \lessgtr
Y'$.
\end{quote}

Pour b), soient $X'' \in \widetilde{X}' \smallsetminus \widetilde{Y}'$ et $Y''
\in \widetilde{Y}' \smallsetminus \widetilde{X}'$. S'ils sont tous deux hors de
$\widetilde{X} \cap \widetilde{Y}$, $X'' \mathrel{|\circ|} Y''$ par définition ;
si par exemple $X'' \in \widetilde{X} \cap \widetilde{Y}$, alors $X''$ et $Y''$
sont deux strates distinctes de $Y$, et Mag 4 conclut.

\begin{quote}
\textbf{Proposition 2} (marge de la page 9). Si $X \neq Y$ et $X \lessgtr Y$,
alors $X \mathrel{|\circ|} Y$.
\end{quote}

Si aucun des deux n'est strate de l'autre, c'est la définition de $\lessgtr$ ;
si $X < Y$, c'est Mag 4. La marge écrit une équivalence ; la réciproque est
fausse\footnote{La marge énonce « si $X \neq Y$, $X \mathrel{|\circ|} Y
\Leftrightarrow X \lessgtr Y$ » et sa démonstration s'interrompt. Contre-exemple
à $\Rightarrow$ : dans le plan, deux segments $X = [p,q]$ et $Y = [r,s]$ dont
l'extrémité $q$ tombe à l'intérieur de $Y$. Leurs intérieurs sont disjoints,
mais le point $q$, strate de $X$, n'est pas intérieurement disjoint de $Y$ : ils
ne sont pas compatibles. Seule l'implication énoncée ici sert dans la suite.},
et c'est elle qui sert :

\begin{quote}
\textbf{Corollaire 1.} Si $X \lessgtr Y$ et $Z \mathrel{\mathring{\ll}} X$,
alors $Z \ll Y \iff X \ll Y$.

\textbf{Corollaire 2.} Si $Z \mathrel{\mathring{\ll}} X$, $Z
\mathrel{\mathring{\ll}} Y$ et $X \lessgtr Y$, alors $X = Y$.

\textbf{Corollaire 3.} Si $X \lessgtr Y$, alors $X \ll Y \iff X \leq Y$.
\end{quote}

Pour le corollaire 1, il suffit de montrer $\Rightarrow$. Par Mag 2 il existe
$Y' \leq Y$ avec $Z \mathrel{\mathring{\ll}} Y'$ ; on a $X \lessgtr Y'$ par la
proposition $1_{0}$. Si $X \neq Y'$, la proposition 2 donne $X
\mathrel{|\circ|} Y'$, puis Mag 3 donne $Z \mathrel{|\circ|} Z$, contre
l'antiréflexivité. Donc $X = Y' \leq Y$. Le corollaire 2 en résulte ($X \ll Y$
et $Y \ll X$), et le corollaire 3 aussi, en prenant $Z = X$ et $Y'$ comme
ci-dessus\footnote{La transcription porte « $X \mathrel{\mathring{\ll}} Y \iff
X \leq Y$ ». Sous cette forme l'énoncé est faux dès que $X < Y$, puisqu'alors le
plus petit élément de $\widetilde{Y}^{X}$ est $X$ et non $Y$ ; ce qui est vrai
avec l'anneau, c'est que deux éléments compatibles vérifiant $X
\mathrel{\mathring{\ll}} Y$ sont égaux. La version pour les figures, page 10,
porte bien $\ll$.}. Ainsi, \emph{entre éléments compatibles, raffinement et
incidence coïncident} : une cellule contenue dans une cellule compatible en est
une face.

\begin{quote}
\textbf{Proposition 3} (page 10). $\mathring{\ll}$ est transitive : $X
\mathrel{\mathring{\ll}} Y \mathrel{\mathring{\ll}} Z$ entraîne $X
\mathrel{\mathring{\ll}} Z$.
\end{quote}

Par Mag 2, $X \mathrel{\mathring{\ll}} Z' \leq Z$ pour un $Z'$. Si $Z' \neq Z$,
Mag 4 donne $Z' \mathrel{|\circ|} Z$, et Mag 3, appliqué deux fois, $X
\mathrel{|\circ|} Y$ puis $X \mathrel{|\circ|} X$ : absurde. Ainsi
$\mathring{\ll}$ est une relation d'ordre, ce qui deviendra un axiome à la page
29.

Le NB de la page 10 ajoute : si $X \ll Y \ll Z$ et $X \mathrel{\mathring{\ll}}
Z$, alors $Y \mathrel{\mathring{\ll}} Z$, mais pas nécessairement $X
\mathrel{\mathring{\ll}} Y$ --- un point intérieur à un segment $Z$ peut être
l'extrémité d'un sous-segment $Y$. Pour les figures enfin : si $F \lessgtr G$,
alors $F \ll G \iff F \leq G$, de sorte que sur les sous-figures d'une figure,
$\ll$ et $\leq$ coïncident.

\pagerange{11}{13}
\section*{Restriction et recollement des raffinements (pages 11 à 13)}

Les propositions 1 et 2 de la page 11 sont les propriétés d'ensemble ordonné de
$\mathrm{Fig}(\mathcal{M})$ : existence des Inf de familles non vides, des Sup
majorés, et $G \lessgtr \bigcup_{i} F_{i}$ si et seulement si $G \lessgtr F_{i}$
pour tout $i$.

\begin{quote}
\textbf{Proposition 3} (page 11). Soient $F, G$ des figures avec $G \ll F$, et
$F' \leq F$. La partie
\[
G_{F'} = G \cap \mathrm{Omb}(F') = \{Y \in G \mid Y \ll X \text{ pour un } X
\in F'\}
\]
est la plus grande sous-figure de $G$ qui raffine $F'$, et $F' \mapsto G_{F'}$
envoie les sous-figures de $F$ dans celles de $G$ en commutant aux Sup et aux
Inf de familles non vides.
\end{quote}

C'est la \emph{restriction} du raffinement $G$ à la sous-figure $F'$ ; la page
l'écrit $G | F'$ ou $\mathrm{Inf}^{\ll}(F', G)$\footnote{La page demande
aussi $F \lessgtr G$, puis le biffe. Le nom $\mathrm{Inf}^{\ll}$ est le sien :
nous n'établissons pas que $G_{F'}$ soit une borne inférieure pour $\ll$ parmi
toutes les figures, seulement qu'elle est la plus grande sous-figure de $G$ qui
raffine $F'$. La démonstration n'est pas sur la page. $G_{F'}$ est fermée parce
que $Y' \leq Y \ll X$ entraîne $Y' \ll X$. La commutation aux Sup est
immédiate. Pour les Inf : si $Y \in G$ raffine un $X_{i} \in F'_{i}$ pour
chaque $i$, soit $Y_{i} \leq X_{i}$ le plus petit élément de Mag 2, qui est
dans $F'_{i}$ ; on a $Y \mathrel{\mathring{\ll}} Y_{i}$, et les $Y_{i}$, tous
dans la figure $F$, sont compatibles, donc égaux par le corollaire 2 de la page
9 : $Y$ raffine un élément de $\bigcap_{i} F'_{i}$. C'est là que sert
l'hypothèse $G \ll F$, qui garantit que les $X_{i}$ sont dans une même
figure.}. Si $G \leq F$, on a simplement $G_{F'} = G \cap F'$ (corollaire 1,
page 12), par le corollaire 3 de la page 9.

\begin{quote}
\textbf{Proposition 4} (page 12, « recollement des raffinements »). Si $F =
\bigcup_{i} F_{i}$, la suite
\[
\mathrm{Raff}(F) \longrightarrow \prod_{i} \mathrm{Raff}(F_{i})
\rightrightarrows \prod_{i,j} \mathrm{Raff}(F_{i} \cap F_{j}),
\]
où $\mathrm{Raff}(F) = \{G \mid G \ll F\}$ et les flèches sont des
restrictions, est exacte. En particulier $\mathrm{Raff}(F) =
\varprojlim_{X \in F} \mathrm{Raff}(\widetilde{X})$.
\end{quote}

Un raffinement d'une figure est donc la même chose qu'une famille de
raffinements de ses cellules qui se recollent sur les faces communes. C'est une
condition de faisceau pour $\mathrm{Raff}$ sur l'ensemble ordonné
$\mathcal{M}$\footnote{Le rapprochement avec les faisceaux (au sens de SGA 4)
sur un ensemble ordonné muni de sa topologie d'Alexandrov est le nôtre. La page
n'en donne pas de démonstration. L'injectivité est claire : un raffinement est
la réunion de ses restrictions, tout élément raffinant un élément de l'un des
$F_{i}$. Le recollement, lui, demande que la réunion de raffinements compatibles
soit formée d'éléments deux à deux compatibles, ce que nous ne vérifions pas ;
la page 19 dira que, dans un atelier, il faut un axiome de plus.}. Le
corollaire 2 de la page 13 en tire que $\leq$, $\ll$, $\lessgtr$, $\parallel$
entre deux raffinements de $F$ se testent sur les restrictions aux $F_{i}$.

\begin{quote}
\textbf{Proposition 5} (page 12). Soient $F \lessgtr G$, $L = F \cap G$, $F'
\ll F$, $G' \ll G$. Alors $F' \lessgtr G'$ si et seulement si $F'_{L} \lessgtr
G'_{L}$, et de même pour $\parallel$.
\end{quote}

Deux raffinements de figures compatibles se rencontrent bien ou mal selon ce
qui se passe sur la partie commune. La page n'en donne pas de démonstration, et
une flèche en marge dit qu'elle « demande la proposition 4 ». Le cas $L =
\emptyset$, lui, est élémentaire : si $F \parallel G$, alors $F' \parallel
G'$\footnote{Si $X' \ll X \in F$ et $Y' \ll Y \in G$, Mag 2 fournit $X_{0} \leq
X$ et $Y_{0} \leq Y$ avec $X' \mathrel{\mathring{\ll}} X_{0}$, $Y'
\mathrel{\mathring{\ll}} Y_{0}$ ; $X_{0} \in F$, $Y_{0} \in G$, donc $X_{0}
\mathrel{|\circ|} Y_{0}$, et Mag 3 conclut.}.

La page 13 constate que tout ceci est « récapitulatif » et ne fait intervenir
$\mathrel{|\circ|}$ qu'à travers $\lessgtr$, et dresse un programme : 1) les
axiomes des lieux et des ombres ; 2) le lien avec le cas ensembliste ; « mais
surtout » 3) revoir la théorie des subdivisions ; 4) le cas où $\ll$ s'exprime
par $\leq$ et $\mathrel{|\circ|}$, qui demande de donner un sens à « $X^{\circ}
\subset Y^{\circ}$ » et donc de passer d'abord par 3). Seuls 1) et 2) sont
traités dans ce dossier.

\pagerange{14}{17}
\section*{Lieux, et les magasins divisibles, quasi-ensemblistes, ensemblistes (pages 14 à 17)}

Les \emph{lieux} sont les éléments minimaux pour $\ll$ ; ils jouent le rôle des
points\footnote{Ce sont aussi les éléments minimaux à la fois pour $\leq$ et
pour $\mathring{\ll}$, comme les définira la page 31 : un élément minimal pour
$\ll$ l'est pour les deux relations plus fines, et réciproquement, si $y \ll x$
avec $x$ minimal pour $\leq$ et $\mathring{\ll}$, Mag 2 donne $y
\mathrel{\mathring{\ll}} x$, donc $y = x$.}. Pour deux lieux, $x \parallel y$
et $x \mathrel{|\circ|} y$ coïncident. On pose
\[
\mathrm{omb}(X) = \{x \in \mathcal{L} \mid x \ll X\}, \qquad
\mathrm{omb}(X)^{\circ} = \mathrm{omb}(X) \smallsetminus \mathrm{omb}(\partial
X) = \{x \in \mathcal{L} \mid x \mathrel{\mathring{\ll}} X\},
\]
l'ombre et l'ombre intérieure de $X$ sur les lieux. Par Mag 3, $X
\mathrel{|\circ|} Y$ entraîne que les lieux intérieurs à $X$ sont disjoints de
ceux intérieurs à $Y$. Deux axiomes possibles :

\begin{quote}
\textbf{Mag L 1.} Pour tout $X$, $\mathrm{omb}(X)^{\circ} \neq \emptyset$.

\textbf{Mag L 2.} Si $x \parallel y$ pour tous $x \in \mathrm{omb}(X)^{\circ}$,
$y \in \mathrm{omb}(Y)^{\circ}$, alors $X \mathrel{|\circ|} Y$.
\end{quote}

Mag L 2 entraîne Mag L 1, par antiréflexivité ; avec Mag 3, il dit que la
disjonction intérieure se lit sur les lieux. « Il semble que chaque fois qu'on a
Mag L 1, on a aussi Mag L 2, qui semble donc l'axiome le plus utile. »

\begin{quote}
\textbf{Scholie} (page 15). Se donner un magasin \emph{strictement divisible}
(satisfaisant Mag L 2) revient à se donner $\leq$, $\ll$ satisfaisant Mag 1,
Mag 2 et Mag L 1, et une relation $\parallel$ \emph{sur les lieux} symétrique et
antiréflexive, telle que si $Y \neq Z$ sont deux strates d'un même $X$, tout
lieu intérieur à $Y$ soit disjoint de tout lieu intérieur à $Z$. On définit
alors $X \mathrel{|\circ|} Y$ par la condition de Mag L 2.
\end{quote}

On vérifie Mag 3 par la transitivité de $\mathring{\ll}$ (page 10), et Mag 4 par
l'hypothèse sur les strates ; la réciproque demande $\mathrm{omb}(X)^{\circ}
\neq \emptyset$, comme le dit la marge.

Le cas \emph{quasi-ensembliste} est celui où deux lieux distincts sont toujours
disjoints. Il ne reste alors que $(\leq, \ll)$ soumis à Mag 1, Mag 2 et Mag L
1, la disjonction intérieure étant $\mathrm{omb}(X)^{\circ} \cap
\mathrm{omb}(Y)^{\circ} = \emptyset$\footnote{Mag 4 n'est pas immédiat : si $x
\mathrel{\mathring{\ll}} Y$ et $x \mathrel{\mathring{\ll}} Z$ avec $Y, Z \leq
X$, le plus petit élément de $\widetilde{X}^{x}$ que fournit Mag 2 est à la
fois $Y$ et $Z$, donc $Y = Z$. Seul Mag 2 sert.}. La page 16 renvoie au
chapitre VII, page 57, pour un exemple montrant qu'un tel magasin ne se plonge
pas forcément dans les figures ensemblistes de ses lieux. Elle note pourtant
que $F' \mapsto \mathrm{mulomb}(F')$ identifie les sous-figures de $F$ à celles
de la figure ensembliste formée des ombres de ses éléments.

Le magasin est \emph{ensembliste} quand, de plus, les relations entre figures
se lisent sur ces ombres\footnote{Page 17 : $\mathrm{mulomb}(F) \ll
\mathrm{mulomb}(G) \Rightarrow F \ll G$ et $\mathrm{mulomb}(F) \lessgtr
\mathrm{mulomb}(G) \Rightarrow F \lessgtr G$, et une condition a) que la page
annonce au bas de la page 16 et n'écrit pas.}, et l'on a alors :

\begin{quote}
\textbf{Scholie} (page 17). Un magasin ensembliste est la donnée d'un ensemble
$L$ (« lieux » ou « points ») et d'une partie $\mathcal{M}$ de l'ensemble
$\mathrm{Figél}(L)$ des figures ensemblistes élémentaires de $L$ telle que :
\end{quote}
\begin{enumerate}
\item[Magens 1.] si $F \in \mathcal{M}$ et $X \in F$, la figure $\{Y \in F \mid
Y \subset X\}$ est dans $\mathcal{M}$ ;
\item[Magens 2.] pour tout $x \in L$, $\{\{x\}\} \in \mathcal{M}$.
\end{enumerate}


Une figure ensembliste élémentaire est une famille admissible de parties de
$L$ ayant un plus grand élément, au sens du chapitre V\footnote{Dossier 156-5,
pages 1 à 5 : une famille de parties est admissible si tout point de sa
réunion est dans un plus petit membre, et ses membres sont alors les strates
fermées d'une topologie d'Alexandrov.}. « À peine plus compliqué que la
définition d'une “topologie” ! », dit la marge, et la page conclut qu'il « vaut
décidément mieux partir du magasin », quitte à introduire ensuite des ateliers
plus généraux que l'atelier canonique.

\pagerange{18}{23}
\section*{Ateliers (pages 18 à 23)}

\subsection*{L'atelier comme partie de $\mathrm{Fig}(\mathcal{M})$}

Le 28 juin (« 27 juin (ou plutôt, 28 juin ! (on ne sait plus …) »), page 18 :

\begin{quote}
\textbf{Définition.} Un \emph{antifiltre} d'un ensemble ordonné est une partie
fermée vers le bas et stable par les bornes supérieures de deux éléments
lorsqu'elles existent. Un \emph{atelier} de $\mathcal{M}$ est un antifiltre
$\mathfrak{F}$ de $\mathrm{Fig}(\mathcal{M})$ contenant les figures
élémentaires $\widetilde{X}$ et la figure vide.
\end{quote}

L'antifiltre est ce qu'on appelle aujourd'hui un idéal de l'ensemble
ordonné\footnote{La notion duale de celle de filtre, d'où le nom. Le mot
« idéal » d'un ensemble ordonné, ou d'un treillis, est postérieur à
l'usage qu'en fait la théorie des anneaux et courant dans les années 1930
(Birkhoff, Stone) ; le rapprochement est le nôtre.}. La condition sur la figure
vide résulte des autres si $\mathcal{M} \neq \emptyset$ ; un magasin vide n'a
que l'atelier $\{\emptyset\}$. Les axiomes « Mag at » de la marge disent la même
chose en termes de parties de $\mathcal{M}$ : stabilité par sous-figure, par
réunion de deux figures compatibles, présence des $\widetilde{X}$, non-vacuité.

Les relations $\leq$, $\ll$, $\lessgtr$, $\parallel$ passent à
$\mathfrak{F}$, et les propositions des pages 9 à 13 restent vraies, sauf le
recollement des raffinements, qui demande un axiome de plus :
\begin{quote}
\textbf{At Fil.} Si $(G_{\alpha})$ est une famille filtrante croissante de
raffinements de $F \in \mathfrak{F}$ dont la restriction à chaque $X \in F$ est
stationnaire, $\bigcup_{\alpha} G_{\alpha}$ est dans $\mathfrak{F}$.
\end{quote}

\subsection*{L'atelier comme ensemble ordonné}

On peut aussi partir de l'atelier seul (pages 19 et 20), un ensemble
$\mathfrak{F}$ muni de $\leq$, $\ll$, $\mathrel{|\circ|}$ :
\begin{enumerate}
\item[At 1.] Les Sup majorés existent, et $\mathfrak{F}$ a un plus petit
élément.
\item[At 2.] Appelant \emph{élémentaires} les $X$ tels que $X = \mathrm{Sup}\,
F_{i}$ entraîne $X = F_{i}$ pour un $i$, et $\mathcal{M}$ leur ensemble, on a
$F \leq G$ si et seulement si toute strate élémentaire de $F$ est strate de $G$
--- « toute figure est le Sup de ses strates ».
\item[At 3.] $\{F, G\}$ est majoré si et seulement si leurs strates
élémentaires sont deux à deux compatibles.
\item[At 4.] $F \ll G$ si et seulement si toute strate élémentaire de $F$
raffine une strate élémentaire de $G$.
\item[At 5.] $\mathrel{|\circ|}$ ne lie que des éléments de $\mathcal{M}$, et
la compatibilité de deux d'entre eux s'exprime par $\mathrel{|\circ|}$ comme à
la page 7.
\end{enumerate}
avec Mag 1 à Mag 4 sur $\mathcal{M}$\footnote{Les éléments élémentaires sont
ce qu'on nomme aujourd'hui les éléments complètement sup-irréductibles, et At 2
dit que tout élément est la borne supérieure de ceux qui sont au-dessous de lui
--- l'analogue, pour un ensemble ordonné qui n'est pas un treillis, de la
représentation de Birkhoff (1937) d'un treillis distributif fini par ses
sup-irréductibles. Le rapprochement est le nôtre.}.

\begin{quote}
\textbf{Scholie} (page 20). Les deux points de vue --- $(\mathcal{M}, \leq,
\ll, \mathrel{|\circ|}, \mathfrak{F} \subset \mathrm{Fig}(\mathcal{M}))$ soumis
à Mag 1 à Mag 4 et Mag at, et $(\mathfrak{F}, \leq, \ll, \mathrel{|\circ|})$
soumis à At 1 à At 5 et Mag 1 à Mag 4 --- donnent des groupoïdes équivalents.
\end{quote}

« Le point de vue “magasin” est plus pratique, visiblement, pour l'économie de
l'écriture des axiomes de base. Le point de vue “atelier” me semble plus
commode pour la pratique du calcul et des raisonnements géométriques, une fois
l'axiomatique acquise. »

\subsection*{Le point de vue mixte (pages 21 à 23)}

Une troisième « prise de vue synthétique » se donne deux ensembles,
$\mathcal{M}$ et $\mathfrak{F}$, et une relation d'incidence $X \lhd F$ entre
eux, avec $\ll$ et $\mathrel{|\circ|}$ sur $\mathcal{M}$. Les axiomes At M 1 à
At M 5 demandent : que $F \mapsto \widetilde{F} = \{X \mid X \lhd F\}$ soit
injective et sépare les éléments de $\mathcal{M}$ ; que pour chaque $X$ il y ait
une plus petite figure $F_{X}$ incidente à $X$, d'où l'ordre $X \leq Y \iff
F_{X} \leq F_{Y}$ ; que $\widetilde{\mathfrak{F}}$ soit stable par réunion de
sous-familles majorées et contienne $\emptyset$ ; que la compatibilité des
figures se teste sur leurs éléments, et celle des éléments par
$\mathrel{|\circ|}$. La page 22 note que toute partie fermée d'un
$\widetilde{F}$ est encore dans $\widetilde{\mathfrak{F}}$, ce qui est
équivalent à At M 3 a), laissé sous sa forme faible pour l'énoncer comme une
existence de Sup majorés.

Bilan, page 23 : partant de $\mathfrak{F}$, ou de $\mathfrak{F}$ et
$\mathcal{M}$, il faut cinq axiomes en plus de Mag 1 à Mag 4 ; partant de
$(\mathcal{M}, \leq, \ll, \mathrel{|\circ|})$, un seul suffit, sur le choix de
$\mathfrak{F}$. Les ateliers de ces pages, il propose de les appeler
« spécieux ».

\pagerange{23}{28}
\section*{La variante polyidéale (pages 23 à 28)}

« Il faut maintenant ouvrir une parenthèse sur les ateliers “polyidéaux”, qui
représentent une variante très légèrement différente. »

\begin{quote}
\textbf{Définitions} (pages 23 à 25). Le magasin est \emph{prépolyidéal} si
deux éléments majorés de $\mathcal{M}$ ou bien n'ont pas de minorant commun, ou
bien ont une borne inférieure. On pose $X \ll_{\mathrm{pol}} Y$ si $X \ll Y$ et
si, pour tout $Y' \leq Y$, $\widetilde{X} \cap \mathrm{Omb}(Y')$ est vide ou a
un plus grand élément. Le magasin est \emph{polyidéal} si $\ll$ et
$\ll_{\mathrm{pol}}$ coïncident.
\end{quote}

La page écrit « ont un Inf »\footnote{La première rédaction de la définition
est raturée et s'enchevêtre avec la seconde. Prise à la lettre (« deux éléments
majorés ont un Inf »), la définition exclurait deux sommets distincts d'un même
segment, qui n'ont aucun minorant ; et elle ne serait plus équivalente à $X \leq
Y \Rightarrow X \ll_{\mathrm{pol}} Y$, que la page 24 lui substitue, où l'on
demande « vide ou a un plus grand élément ». Nous lisons la définition dans ce
sens.} ; avec cette lecture, $\mathcal{M}$ est prépolyidéal si et seulement si
$X \leq Y$ entraîne $X \ll_{\mathrm{pol}} Y$ (page 24). Les pages 24 et 25
vérifient alors que $(\mathcal{M}, \leq, \ll_{\mathrm{pol}}, \mathrel{|\circ|})$
est encore un magasin, le \emph{magasin polyidéal associé} : la transitivité de
$\ll_{\mathrm{pol}}$ par un argument de sous-figures induites, Mag 2 parce que
$X \ll Y' \leq Y$ entraîne déjà $X \ll_{\mathrm{pol}} Y'$, et Mag 3, Mag 4 sans
changement.

Il faut ensuite distinguer deux compatibilités : $X \lessgtr_{\mathrm{pol}} Y$
si $X \lessgtr Y$ et si deux strates quelconques $X' \leq X$, $Y' \leq Y$ ont une
intersection $\widetilde{X}' \cap \widetilde{Y}'$ vide ou ayant un plus grand
élément. Les \emph{figures polyidéales}, $\mathrm{Figpol}(\mathcal{M}) \subset
\mathrm{Fig}(\mathcal{M})$, sont celles dont les éléments sont deux à deux
$\lessgtr_{\mathrm{pol}}$, et :

\emph{même si $\mathcal{M}$ est polyidéal, l'inclusion est stricte en
général.} Le croquis de la page 26 est un cercle portant deux points : deux
arêtes ayant les mêmes deux extrémités forment une figure, mais leurs
strates communes, les deux sommets, n'ont pas de plus grand élément. C'est
exactement ce qui sépare un complexe polyédral, où deux cellules se coupent
selon une face commune, d'un complexe CW régulier, où elles peuvent se couper
selon plusieurs\footnote{Le cercle découpé en deux arêtes et deux sommets est
l'exemple le plus simple de complexe CW régulier qui n'est pas un complexe
simplicial ni polyédral. La marge donne des analogues en dimension 2 (deux
triangles), cette fois pour $X \ll Y$ sans $X \ll_{\mathrm{pol}} Y$. Les noms
sont les nôtres.}.

Pour $F, G \in \mathrm{Figpol}(\mathcal{M})$ (page 26), les conditions
suivantes sont équivalentes : $\{F,G\}$ majoré, $\mathrm{Sup}(F,G)$ existe,
$F \cup G \in \mathrm{Figpol}(\mathcal{M})$, et $X \lessgtr_{\mathrm{pol}} Y$
pour tous $X \in F$, $Y \in G$. Un scholie de la même page annonce que
les propriétés algébriques de $\mathrm{Figpol}(\mathcal{M})$, avec $\leq$,
$\ll_{\mathrm{pol}}$ et $\lessgtr_{\mathrm{pol}}$, seront « essentiellement les
mêmes » que celles de $\mathrm{Fig}(\mathcal{M})$, la seule différence étant
l'expression de $\lessgtr_{\mathrm{pol}}$ par $\mathrel{|\circ|}$. Un \emph{atelier polyidéal} se définit comme
l'atelier spécieux, avec $\mathrm{Figpol}$ au lieu de $\mathrm{Fig}$ (axiomes
« Mag at pol », page 27). Un atelier spécieux est polyidéal si et seulement si
le magasin l'est et si $X \lessgtr Y$ entraîne $X \lessgtr_{\mathrm{pol}} Y$,
« ce qui est beaucoup plus exceptionnel ».

Conclusion de la page 27 : ce sont « des notions très voisines, mais
disjointes » ; les notions spécieuses sont commodes pour travailler avec des
figures aussi générales que possible, les polyidéales pour découper en
morceaux « aussi standard que possible ». La marge en donne la raison : on a
pris $\ll$ comme relation primitive, et non $\lessgtr$, « il y a un certain
arbitraire ». Dans le magasin, la différence se réduit à un axiome (page 28) :
\begin{quote}
\textbf{Mag pol.} Si $X \ll Y$ et $Y' \leq Y$, $\widetilde{X} \cap
\mathrm{Omb}(Y')$ est vide ou a un plus grand élément.
\end{quote}

\pagerange{28}{32}
\section*{L'incidence intérieure comme donnée primitive (pages 28 à 32)}

\subsection*{Mag 1' à Mag 4'}

La relation $\ll$ se retrouve à partir de $\mathring{\ll}$ :
\[
(*) \qquad X \ll Y \iff X \mathrel{\mathring{\ll}} Y' \leq Y \text{ pour un }
Y'.
\]
Les pages 28 et 29 cherchent quand une structure $(\mathcal{M}, \leq,
\mathring{\ll}, \mathrel{|\circ|})$ provient ainsi d'un magasin dont
$\mathring{\ll}$ soit l'incidence intérieure, et aboutissent page 30 à :

\begin{quote}
\textbf{Définition} (pages 29 et 30). Un magasin est un ensemble $\mathcal{M}$
muni de $\leq$, $\mathring{\ll}$, $\mathrel{|\circ|}$ tels que :
\end{quote}
\begin{enumerate}
\item[Mag 1'.] $\leq$ et $\mathring{\ll}$ sont des relations d'ordre ;
$\mathrel{|\circ|}$ est symétrique et antiréflexive.
\item[Mag 2'.] (Entre $\leq$ et $\mathring{\ll}$.) a) Si $X
\mathrel{\mathring{\ll}} Y' \leq Y$, $Y'$ est déterminé par $X$ et $Y$. b) Si
$X \mathrel{\mathring{\ll}} Y$ et $X' \leq X$, il existe $Y' \leq Y$ avec $X'
\mathrel{\mathring{\ll}} Y'$.
\item[Mag 3'.] (Entre $\mathring{\ll}$ et $\mathrel{|\circ|}$.) Si $X
\mathrel{|\circ|} Y$, $X' \mathrel{\mathring{\ll}} X$, $Y'
\mathrel{\mathring{\ll}} Y$, alors $X' \mathrel{|\circ|} Y'$.
\item[Mag 4'.] (Entre $\mathrel{|\circ|}$ et $\leq$.) Comme Mag 4.
\end{enumerate}


La page 30 récrit Mag 4' sous la forme « si $X \leq Y$ et $X \neq Y$, alors $X
\mathrel{|\circ|} Y$ », précédée d'un « non » douteux que nous
écartons\footnote{Avec « non », l'axiome contredirait Mag 4, Magor 4 (page 3) et
la page 29, qui dit « Mag 3', Mag 4' comme Mag 3, Mag 4 ». Sans lui, la forme
de la page 30, restreinte aux éléments comparables, est \emph{plus faible} que
Mag 4 : sur $\mathcal{M} = \{X, Y, Z\}$ avec $Y, Z < X$ incomparables,
$\mathring{\ll}$ l'égalité, et $\mathrel{|\circ|}$ reliant seulement $X$ à $Y$
et à $Z$, Mag 1' à 3' et la forme restreinte sont satisfaits, mais pas Mag 4.
Nous gardons Mag 4.}. La page 30 dit pourquoi $\mathring{\ll}$ est « plus
primaire » que $\ll$ et « plus intuitive » : elle est disjointe de $\leq$
($X \leq Y$ et $X \mathrel{\mathring{\ll}} Y$ entraînent $X = Y$), et elle est
celle qui figure naturellement dans Mag 3 ; et chaque axiome lie désormais deux
relations et deux seulement. « Décidément, c'est la forme parfaite de
définition ! »

\subsection*{Trois semaines de choix}

Suit, page 30, le « rappel des principales phases parcourues » pour le choix
des données de base, depuis le chapitre IV, page 52, le 13
juin\footnote{Dossier 156-4. La page dit « depuis quinze semaines », le mot
« semaines » lu avec doute ; du 13 juin au 28 juin il y a quinze jours.} :
\begin{enumerate}
\item[--] des figures, avec $\leq$ et une relation de subdivision ;
\item[--] des figures, avec $\leq$ et $\ll$ ;
\item[--] $\mathcal{M}$, $\leq$, $\ll$, $\lessgtr$ (« retour vers les
origines ») ;
\item[--] $\mathcal{M}$, $\leq$, $\ll$, $\mathrel{|\circ|}$, « l'innovation
d'avant-hier faisant tomber tous les fouillis des axiomes de compatibilité --
disjonction » ;
\item[--] $\mathcal{M}$, $\leq$, $\mathring{\ll}$, $\mathrel{|\circ|}$,
« aujourd'hui, l'épanouissement intégral ! ».
\end{enumerate}
Seule des trois relations, $\leq$ passe aux figures ; d'où (page 31) l'idée de
se donner $(\mathcal{M}, \mathfrak{F} \mid \lhd, \mathring{\ll},
\mathrel{|\circ|})$, l'ordre $\leq$ sortant de l'incidence $\lhd$.

\subsection*{Les cas particuliers, dans la nouvelle forme}

\emph{Magasins strictement divisibles} (page 31). Les lieux sont les éléments
minimaux pour $\leq$ et pour $\mathring{\ll}$ ; les données sont $(\mathcal{M},
\leq, \mathring{\ll}, \parallel)$ avec $\parallel$ symétrique et antiréflexive
sur $\mathcal{L}$, Mag div 2 étant Mag 2', et Mag div 3 : si $X < Y$, tout lieu
intérieur à $X$ est disjoint de tout lieu intérieur à $Y$. On retrouve $X
\mathrel{|\circ|} Y$ si et seulement si tous leurs lieux intérieurs sont
disjoints\footnote{Mag div 3 est ici, comme Mag 4' à la page 30, énoncé pour
des éléments comparables ; la forme de la page 15 (deux strates distinctes d'un
même $X$) est celle qui donne Mag 4.}.

\emph{Magasins quasi-ensemblistes} (page 32). Les données sont $(\mathcal{M},
\leq, \mathring{\ll})$ avec Mag quens 1 ($\leq$, $\mathring{\ll}$ des ordres),
Mag quens 2 (Mag 2') et Mag quens 3 (tout $X$ a un lieu $x
\mathrel{\mathring{\ll}} X$)\footnote{La page porte « $x
\mathrel{\mathring{\ll}} \mathcal{M}$ » ; il faut lire $X$.}, la disjonction
intérieure étant $\mathrm{Omb}^{\circ}(X) \cap \mathrm{Omb}^{\circ}(Y) \cap
\mathcal{L} = \emptyset$.

\emph{Magasins ensemblistes}. La structure de la page 17, « à laquelle je n'ai
rien à ajouter », avec $X \mathrel{|\circ|} Y \iff X^{\circ} \cap Y^{\circ} =
\emptyset$, $X^{\circ} = X \smallsetminus \partial X$.

La page 32 se termine par un tableau des systèmes d'axiomes, que voici avec nos
numéros de page\footnote{Les renvois de la page sont dans sa pagination ; deux
d'entre eux (ses pages 30 et 31, nos pages 38 et 39, peut-être « 31, 33 »)
pointent vers la suite, de sorte que le tableau a été complété plus tard.} :
\begin{enumerate}
\item[--] magasins : Mag 1 à Mag 4 (page 30), Mag L 1, L 2 (pages 14 et 15),
Mag at (page 18) ;
\item[--] ateliers : At 1 à At 5 et Mag 1 à Mag 4 (pages 19 et 20) ;
\item[--] mixtes : At M 1 à At M 5 et Mag 1 à Mag 4 (pages 21 et 22) ;
\item[--] magasins strictement divisibles : Mag div 1 à 3 (page 31) ;
quasi-ensemblistes : Mag quens 1 à 3 (page 32) ; ensemblistes : Magens 1, 2
(page 17) ;
\item[--] variante polyidéale : Mag 1 à Mag 4, Mag pol, Mag at pol (pages 23 à
27).
\end{enumerate}

\pagerange{33}{35}
\section*{Magasins triviaux et figures indexées (pages 33 à 35)}

Un magasin est \emph{trivial} si $\mathring{\ll}$ est l'égalité et
$\mathrel{|\circ|}$ la différence : il ne reste qu'un ensemble ordonné
$(\mathcal{M}, \leq)$, et tout ensemble ordonné en fournit un --- ce sont les
« maquettes » de la table de la page 2. On y trouve $\ll = \leq$, $\lessgtr$
toujours vraie, $X \parallel Y$ si et seulement si $\{X, Y\}$ n'est pas minoré ;
les lieux sont les éléments minimaux, et $\mathrm{omb}(X)^{\circ}$ vaut $\{X\}$
si $X$ est minimal, $\emptyset$ sinon. Un magasin trivial ne satisfait donc Mag
L 1 que si l'ordre est discret.

\begin{quote}
\textbf{Proposition} (pages 33 et 34). Soit $I$ un ensemble ordonné. Une
application $\varphi : I \to \mathcal{M}$, $i \mapsto X_{i}$, définit une
\emph{figure indexée par $I$} --- une figure $F$ avec un isomorphisme ordonné $I
\simeq F$ --- si et seulement si : a) $\varphi$ est croissante et injective ;
b) pour tout $i$, $\varphi$ induit une bijection $I_{\leq i} \to
\widetilde{X}_{i}$ ; c) si $i \neq j$, $X_{i} \mathrel{|\circ|} X_{j}$.
\end{quote}

La condition c) dit, remarque-t-il, que $\varphi$ envoie la structure de
magasin triviale de $I$ dans $\mathcal{M}$. b) rend $\varphi(I)$ fermée ; pour
la compatibilité de $X_{i}$ et $X_{j}$, deux strates $X' = X_{i'}$, $Y' =
X_{j'}$ hors de la partie commune sont distinctes, donc $i' \neq j'$, et c)
conclut (page 35)\footnote{Que $\varphi$ soit un isomorphisme sur son image, et
pas seulement une bijection croissante, résulte de b) et de l'injectivité : si
$X_{i} \leq X_{j}$, alors $X_{i} = X_{i'}$ pour un $i' \leq j$, et $i = i'$.}.

La page 35 ouvre ensuite une section « Variantes polyidéales » de quelques
lignes : la condition polyidéale sur une figure $F$ se lit sur le magasin
$\widetilde{F}$, et dit que c'est un ensemble ordonné polyidéal.

\pagerange{35}{40}
\section*{Morphismes de magasins (pages 35 à 40)}

\begin{quote}
\textbf{Définition} (pages 35 et 36). Un \emph{morphisme} $\varphi : \mathcal{M}
\to \mathcal{M}'$ est une application telle que a) $X \leq Y \Rightarrow
\varphi(X) \leq \varphi(Y)$ ; b) $X \mathrel{\mathring{\ll}} Y \Rightarrow
\varphi(X) \mathrel{\mathring{\ll}} \varphi(Y)$ ; c) $X \mathrel{|\circ|} Y
\Rightarrow \varphi(X) \mathrel{|\circ|} \varphi(Y)$ ; et, « il semble qu'on
ait besoin aussi », d) pour tout $X$, $\varphi$ induit une bijection
$\widetilde{X} \to \widetilde{\varphi(X)}$.
\end{quote}

La marge se demande si c'est « trop exigeant », et suggère qu'on pourrait se
contenter de $X \lessgtr Y \Rightarrow \varphi(X) \lessgtr \varphi(Y)$. Une
autre marge, en regard de d), datée 11.7.86\footnote{Ou 1.7.86 ; la lecture
11.7 placerait cette note pendant le chapitre IX. Elle est presque
illisible.}, met d) en question.

\begin{quote}
\textbf{Corollaires} (pages 36 et 37). Sous a), c), d) : 1) l'image d'une
figure $F$ est une figure, et $\varphi|_{F}$ est un isomorphisme ordonné sur
elle ; 2) $\varphi$ préserve $\lessgtr$ et $\parallel$ ; 3) sur une figure,
$\leq$ et $\parallel$ sont reflétés ; 4) $\varphi$ commute à la réunion et à
l'intersection de deux figures compatibles ; 5) $\varphi$ préserve le
raffinement des figures, et 6) la restriction des raffinements.
\end{quote}

Le point du corollaire 1 est l'injectivité sur $F$ : deux éléments distincts de
$F$ sont compatibles, donc intérieurement disjoints (proposition 2 de la page
9), et c) les sépare ; la proposition des pages 33 et 34 fait le reste.

« Je n'ai pas l'impression d'avoir bien saisi le rôle exact de l'hypothèse c),
qui semble très forte » (page 37) : pour une image inverse par une application
d'ensembles, elle paraît exiger la surjectivité. La marge propose plutôt
l'implication réciproque, $\varphi(X) \mathrel{|\circ|} \varphi(Y) \Rightarrow X
\mathrel{|\circ|} Y$.

\subsection*{Exemple 1, et les magasins où $\mathrel{|\circ|}$ vient de
$\mathring{\ll}$}

L'application $X \mapsto \mathrm{Multombs}(X)$ de $\mathcal{M}$ dans les
figures ensemblistes élémentaires de $\mathcal{M}$\footnote{Lecture incertaine
des abréviations « Figél » et « Multombs » ; nous les prenons pour
« figures élémentaires » et « multi-ombres », la figure ensembliste formée des
ombres $\mathrm{Omb}(Y)$, $Y \leq X$, de support $\mathrm{Omb}(X)$ et
d'intérieur $\mathrm{Omb}^{\circ}(X)$.} respecte $\leq$, $\ll$, $\lessgtr$, et
$X \mathrel{|\circ|} Y$ entraîne que $X$ et $Y$ n'ont aucun
$\mathring{\ll}$-minorant commun (sinon Mag 3 donnerait $Z \mathrel{|\circ|} Z$)
(page 38). La réciproque est fausse en général : deux lieux distincts n'ont
jamais de minorant commun, et ne sont pas pour autant disjoints. Qu'elle soit
vraie est une hypothèse « de nature quasi-ensembliste, mais sans référence
particulière aux lieux ».

\begin{quote}
\textbf{Proposition} (page 39). Les structures $(\mathcal{M}, \leq,
\mathring{\ll})$ qui deviennent des magasins quand on pose
\[
X \mathrel{|\circ|} Y \iff \{X, Y\} \text{ n'a pas de minorant pour }
\mathring{\ll}
\]
sont exactement celles qui satisfont Mag 1' et Mag 2'.
\end{quote}

Mag 3' résulte de la transitivité de $\mathring{\ll}$, et Mag 4 de l'unicité
de Mag 2' a)\footnote{Si $W \mathrel{\mathring{\ll}} Y$ et $W
\mathrel{\mathring{\ll}} Z$ avec $Y, Z \leq X$, Mag 2' a) appliqué à $W$ et $X$
donne $Y = Z$. C'est la remarque, biffée, de la page 39 : « Mag 4' résulte de
Mag 2' a) ! ».}. « C'est pratiquement aussi simple que le cas ensembliste ! »
Ce sont les \emph{magasins fidèles} de la table de la page 2. D'où une
hiérarchie : magasins généraux $(\leq, \mathring{\ll}, \mathrel{|\circ|})$ ;
magasins fidèles $(\leq, \mathring{\ll})$ ; quasi-ensemblistes ; ensemblistes.
Un quasi-ensembliste est fidèle : un minorant commun à $X$ et $Y$ porte un lieu
intérieur, commun aux deux.

\emph{Exemple 2} (page 40). Sous Mag L 1, $X \mapsto \mathrm{multombs}(X)$, à
valeurs dans les figures ensemblistes des lieux, respecte aussi
$\mathrel{|\circ|}$. Ces exemples montrent que la notion « n'est pas trop
idiote ». \emph{Exemple 3} : une application surjective $f : E' \to E$ induit
un morphisme injectif $\mathrm{Figél}(E) \to \mathrm{Figél}(E')$ ; l'exemple
est repris page 54.

\pagerange{41}{44}
\section*{Fonction support (pages 41 à 44)}

« J'en viens maintenant à la question d'une fonction genre “support” »
$X \mapsto |X|$, de $\mathcal{M}$ dans les parties d'un ensemble $P$ de
« points ». Avec $|\partial X| = \bigcup_{Y < X} |Y|$ et $|X|^{\circ} = |X|
\smallsetminus |\partial X|$, le minimum demandé est\footnote{Comparer, dit la
page, au chapitre VII, page 73 ; le sigle, lu « CF » sur la page, est sans doute
« GF ».} :
\begin{enumerate}
\item[(i)] $X \ll Y \Rightarrow |X| \subset |Y|$ ;
\item[(ii)] $|X|^{\circ} \neq \emptyset$ ;
\item[(iii)] $X \mathrel{|\circ|} Y \iff |X|^{\circ} \cap |Y|^{\circ} =
\emptyset$ ;
\item[(iv)] (normalisation) les $|X|$ recouvrent $P$ et en séparent les
points.
\end{enumerate}
Pour que les $|X|^{\circ}$ séparent encore les points, il faut renforcer (ii) :
\begin{enumerate}
\item[(ii bis)] pour tout $X$, la famille $(|Y|)_{Y \leq X}$ est une figure
ensembliste élémentaire indexée par $\widetilde{X}$, de support $|X|$.
\end{enumerate}
Il en déduit (pages 42 et 43) que l'intérieur $|Y|^{\circ}$ calculé dans cette
famille est le même que celui défini plus haut, et
\begin{enumerate}
\item[(ii ter)] $|X| = \bigcup_{Y \leq X} |Y|^{\circ}$, réunion disjointe de
parties non vides ;
\end{enumerate}
puis, pour $Y, Z \leq X$, que $|Y| \subset |Z|$ entraîne $Y \leq Z$ : avec $L =
\widetilde{Y} \cap \widetilde{Z}$, on a $|Y| \cap |Z| = \bigcup_{W \in L}
|W|^{\circ}$, et $|Y| \subset |Z|$ force $\bigcup_{W \in \widetilde{Y}
\smallsetminus L} |W|^{\circ} = \emptyset$, donc $\widetilde{Y} = L$, les
intérieurs étant non vides\footnote{La page écrit « $\neq$ » là où l'argument
demande « $= \emptyset$ ». Dans la seconde réunion elle écrit $|T|^{\circ}$ pour
$|W|^{\circ}$.}. Inversement, sous (ii ter), tout point de $|X|$ est dans un
plus petit $|Y|$, $Y \leq X$ (page 44) --- « dans tous ces cas, la relation
$\ll$ n'a pas joué de rôle ». Un lemme général commencé au bas de la page 44
est barré, et repris à la page suivante.

\pagerange{45}{53}
\section*{Familles indexées de parties (pages 45 à 53)}

« Je vais reprendre ici la proposition 1, page 1 » du chapitre V\footnote{Dossier
156-5, page 1 ; le sigle est encore lu « CF V ».}, dans un cadre indexé. Soient
$I$ un ensemble ordonné, $X$ un ensemble, et $\Phi = (X_{i})_{i \in I}$ une
famille de parties de $X$ telle que $i \leq j$ entraîne $X_{i} \subset
X_{j}$\footnote{Cette hypothèse est en marge de la page 45 ; la page 47 se
demande « dans l'énoncé de la prop., on n'a pas supposé que $i \leq j
\Rightarrow X_{i} \subset X_{j}$ ?! ». Elle est nécessaire : sans elle, la
condition b) de (iii) ci-dessous n'équivaut plus à 3°.}. On pose
\[
\partial X_{i} = \bigcup_{j < i} X_{j}, \qquad X_{i}^{\circ} = X_{i}
\smallsetminus \partial X_{i}, \qquad |\Phi| = \bigcup_{i} X_{i}, \qquad
I^{x} = \{i \mid x \in X_{i}\},
\]
de sorte que $x \in X_{i}^{\circ}$ si et seulement si $i$ est minimal dans
$I^{x}$.

\begin{quote}
\textbf{Proposition} (pages 45 à 50). Les conditions suivantes sont
équivalentes :
\end{quote}
\begin{enumerate}
\item[(i)] pour tout $x \in |\Phi|$, $I^{x}$ a un plus petit élément ;
\item[(ii)] pour tout $x \in |\Phi|$, $I^{x}$ a un élément minimal, et $X_{i}
\cap X_{j} = \bigcup_{k \leq i, j} X_{k}$ pour tous $i, j$ ;
\item[(iii)] a) les $X_{i}^{\circ}$ sont deux à deux disjoints ; b) $X_{i} =
\bigcup_{j \leq i} X_{j}^{\circ}$ pour tout $i$ ;
\item[(iv)] avec $I' = \{i \mid X_{i}^{\circ} \neq \emptyset\}$ : $\alpha$) les
$X_{i}^{\circ}$, $i \in I'$, forment une partition de $|\Phi|$ ; $\beta$)
chaque $X_{i}$ est saturé pour cette partition ; $\gamma$) si $X_{j}^{\circ}
\neq \emptyset$ et $X_{j}^{\circ} \subset X_{i}$, alors $j \leq i$.
\end{enumerate}
Les conditions « $I^{x}$ a un élément minimal » et b) sont automatiques si $I$
est fini.


La démonstration décompose (i) en trois conditions sur chaque $I^{x}$ : un
élément minimal existe, il est unique, tout élément le majore ; leurs
traductions globales sont $|\Phi| = \bigcup_{i} X_{i}^{\circ}$, la disjonction
a), et b), la dernière entraînant la première. D'où (i) $\iff$ (iii). Pour (iv),
$\alpha$) et $\beta$) redonnent a) et b) dès qu'on sait que $X_{j}^{\circ}
\subset X_{i}$ seulement si $j \leq i$, c'est-à-dire $\gamma$), et $\gamma$)
résulte de (i) :

\begin{quote}
\textbf{Corollaire 1} (page 47). Sous (i), et si $X_{j}^{\circ} \neq
\emptyset$, on a $j \leq i \iff X_{j} \subset X_{i} \iff X_{j}^{\circ} \subset
X_{i} \iff X_{j}^{\circ} \cap X_{i} \neq \emptyset$.

\textbf{Corollaire 2} (page 48). Si de plus tous les $X_{i}^{\circ}$ sont non
vides, $i \mapsto X_{i}$ est un plongement d'ensembles ordonnés. Les conditions
(i) et « $X_{i}^{\circ} \neq \emptyset$ pour tout $i$ » sont donc nécessaires
et suffisantes pour que $\Phi$ soit une figure ensembliste indexée par $I$.
\end{quote}

Ces corollaires portent sur la page un nom lu « Main »\footnote{« Cor.
\uncertain{Main} 1 » : lecture douteuse.}. Pour (i) $\Rightarrow$ (ii) (page 49),
b) et a) donnent $X_{i} \cap X_{j} = \bigcup_{k \leq i,j} X_{k}^{\circ}$. Pour
(ii) $\Rightarrow$ (i), un lemme : la formule d'intersection entraîne a), car
tout $k \leq i, j$ avec $i \neq j$ est strictement sous l'un des deux ; et elle
entraîne que $X_{j}^{\circ} \cap X_{i} \neq \emptyset$ force $j \leq
i$\footnote{La page 50 conclut que $X_{j}^{\circ} \cap X_{i}$ est contenu dans
son complémentaire, « ce qui contredit l'hypothèse que c'est vide » : il faut
lire « non vide ».}, ce qui, avec l'élément minimal de $I^{x}$, donne b).

Le corollaire 3 (page 50), en partie raturé, caractérise les figures indexées
par la disjonction des intérieurs, leur non-vacuité, et le fait que chaque
$\Phi_{i} = (X_{j})_{j \leq i}$ soit une figure ; la page 51, qui le démontre,
est entièrement barrée.

\subsection*{Deux contre-exemples (pages 52 et 53)}

La question de la marge de la page 46 est de savoir si $\beta$) et $\gamma$)
résultent de $\alpha$). Non :

\emph{$\alpha$ sans $\beta$} (page 52). Soient $A_{0} \supsetneq A_{1}
\supsetneq \cdots$ avec $N = \bigcap_{n} A_{n} \neq \emptyset$, et $A_{\infty}$
avec $A_{\infty} \cap A_{0} = N$ et $A_{\infty} \neq N$ ; on indexe par
l'inclusion, $\infty$ étant isolé. Alors $A_{n}^{\circ} = A_{n} \smallsetminus
A_{n+1}$ et $A_{\infty}^{\circ} = A_{\infty}$ sont non vides, deux à deux
disjoints, de réunion $A_{0} \cup A_{\infty}$ ; mais $A_{n}$ n'est pas réunion
des intérieurs qu'il contient --- il manque $N$ ---, et un point de $N$ n'est
dans aucun plus petit membre. $\alpha$) et $\gamma$) sont vrais, $\beta$) est
faux.

\emph{$\alpha$ et $\beta$ sans $\gamma$} (page 53), « moins idiot », même pour
une indexation qui respecte les inclusions : on ajoute $A = A_{\infty}
\smallsetminus N$. Alors $A_{\infty}^{\circ} = N$, $A^{\circ} = A$, $A_{n}^{\circ}
= A_{n} \smallsetminus A_{n+1}$\footnote{La page écrit $A_{i} - A_{i-1}$ ; la
suite étant décroissante, c'est $A_{i+1}$, comme page 52.} : ils sont non vides,
deux à deux disjoints, de réunion $A_{0} \amalg A$, chaque membre est saturé ;
mais $A_{\infty}^{\circ} = N \subset A_{n}$ sans que $A_{\infty} \subset
A_{n}$.

C'est à ces deux pages que renvoie une marge de la page 3 du chapitre V pour un
contre-exemple correct\footnote{Voir la lecture du dossier 156-5, page 3, où le
contre-exemple biffé est faux.}.

\pagerange{53}{60}
\section*{Images inverses, quotients, sous-magasins (pages 53 à 60)}

La section s'intitule « Homomorphismes de magasins et d'ateliers ». Elle
commence par des exemples.

\emph{Exemple 1 : image inverse} (pages 54 à 57). Une application $f :
\mathcal{L}' \to \mathcal{L}$ définit
\[
f^{*}(\Phi) = \{ f^{-1}(A) \mid A \in \Phi,\ f^{-1}(A^{\circ}) \neq \emptyset
\}
\]
sur les figures ensemblistes de $\mathcal{L}$. Avec $\mathcal{L}_{0} =
f(\mathcal{L}')$ et $i : \mathcal{L}_{0} \hookrightarrow \mathcal{L}$, $f^{*}$
se décompose en trois temps : la trace $\Phi \mapsto F_{\mathcal{L}_{0}} \cdot
\Phi$ sur $\mathcal{L}_{0}$, où $F_{\mathcal{L}_{0}} = \{\mathcal{L}_{0}\}$ ;
l'identification des figures de support dans $\mathcal{L}_{0}$ aux figures de
$\mathcal{L}_{0}$ ; et $f_{0}^{*}$, qui pour $f_{0}$ surjective est simplement
$A \mapsto f^{-1}(A)$. Le premier temps est un \emph{passage au quotient} : on
déclare nulle toute multistructure $X$ telle que $X^{\circ} \cap
\mathcal{L}_{0} = \emptyset$, c'est-à-dire $F_{\mathcal{L}_{0}}
\mathrel{|\circ|} X$. En général, conjecture-t-il, un quotient revient à se
donner une partie de $\mathcal{M}$ fermée pour $\mathrel{|\circ|}$ --- un
support, au sens de la page 73 --- dont on néglige les éléments ; négliger $X$,
c'est négliger aussi les $Y \mathrel{\mathring{\ll}} X$, mais pas toutes ses
strates.

L'image directe $i_{*}$ plonge les figures de $\mathcal{L}_{0}$ dans celles de
$\mathcal{L}$, et (page 56)
\[
i^{*} i_{*} = \mathrm{id}, \qquad i_{*} i^{*}(Y) = i_{*}(\mathbf{1}) \cdot Y,
\qquad i_{*}(X') \cdot Y = i_{*}(X' \cdot i^{*}Y),
\]
la dernière étant la \emph{formule de projection}, sous son nom\footnote{C'est
la formule de projection de la théorie des faisceaux et de la cohomologie
étale (SGA 4), dont il est l'auteur ; $\mathbf{1}$ est la figure
$\{\mathcal{L}\}$ à une seule strate, élément unité pour le produit des figures
et plus grand élément pour le raffinement. La page parle d'un « plus petit »
élément $F_{\mathcal{L}}$, lecture incertaine.}. La différence cruciale entre
les deux : $f^{*}$ envoie l'unité sur l'unité et commute (sur les supports) au
complémentaire, $i_{*}$ non.

La difficulté (page 57) : pour $f$ non surjective, $X \mathrel{\mathring{\ll}}
Y$ entraîne bien $f^{*}X \mathrel{\mathring{\ll}} f^{*}Y$, mais pas $X \leq Y
\Rightarrow f^{*}X \leq f^{*}Y$, ni même $f^{*}X \ll f^{*}Y$, puisque $f^{*}Y$
peut être vide et $f^{*}X$ non. Ces applications « inverses » qui comportent un
passage au quotient ne relèvent pas d'une notion « pousse-bouton »\footnote{La
page 57 porte « pour-bonhomme », lu tel quel, et la page 58
« pousse-bouton ».} de morphisme ; en revanche $X \mathrel{|\circ|} Y \Rightarrow
f^{*}X \mathrel{|\circ|} f^{*}Y$ dans tous les cas.

\emph{Exemple 2 : produit par une figure fixe} (page 58). $F \mapsto F \times
F' = \mathrm{pr}_{1}^{*}(F) \cdot \mathrm{pr}_{2}^{*}(F')$, composé de
$\mathrm{pr}_{1}^{*}$ et de la multiplication par une figure fixe : une
multistructure n'y est pas envoyée sur une multistructure.

\emph{Exemple 3 : sous-magasins} (pages 59 et 60). Se restreindre à un type de
multistructures --- connexes, irréductibles, « boules », « simplexes » ---
assez nombreuses pour « résumer » $\mathcal{M}$ : on s'attend à ce que
l'inclusion conserve $\leq$, $\mathrel{|\circ|}$, $\parallel$, mais que
$\mathring{\ll}$ et $\ll$ y deviennent leurs versions polyidéales. C'est un
plongement, sans les difficultés de l'exemple 1. Ce qui conduit à :

\pagerange{61}{66}
\section*{Sous-magasins et homomorphismes d'ateliers (pages 61 à 66)}

\begin{quote}
\textbf{Définition-proposition} (pages 60 et 61). Une partie $\mathcal{M}'$ de
$\mathcal{M}$ est un \emph{sous-magasin} si a) elle est fermée pour $\leq$ ;
b) si $X \mathrel{\mathring{\ll}} Y' \leq Y$ avec $X, Y \in \mathcal{M}'$, alors
$Y' \in \mathcal{M}'$. Munie des relations induites, c'est un magasin, et ses
relations $\ll$, $\lessgtr$, $\parallel$ sont induites par celles de
$\mathcal{M}$\footnote{La condition b) est exactement ce qu'il faut pour que
$\ll$, définie dans $\mathcal{M}'$ par la formule $(*)$ de la page 28, soit la
restriction de celle de $\mathcal{M}$. La page 60 écrit « $\ll_{M'} = \ll_{M}
\mid M$ » : lire $\mid \mathcal{M}'$.}. « Tout marche ! »
\end{quote}

Si $\mathfrak{F}$ est un atelier de $\mathcal{M}$, $\mathfrak{F}' = \{F \in
\mathfrak{F} \mid F \subset \mathcal{M}'\}$ est un atelier de $\mathcal{M}'$,
et l'inclusion commute aux Sup et aux Inf de familles non vides\footnote{À
partir d'ici et jusqu'à la page 67, la page dit « atlas » pour « atelier ».}.
Les propriétés « strictement divisible »\footnote{La transcription lit
« strict. dirigé », avec doute, ici et page 76 ; le renvoi à sa page 23 (notre
page 31, Mag div) impose « strictement divisible ».}, quasi-ensembliste,
ensembliste passent à $\mathcal{M}'$ dès que $\mathcal{M}'$ contient tous les
lieux ; la propriété « $\mathrel{|\circ|}$ définie par $\mathring{\ll}$ » ne
passe que si $\mathcal{M}'$ est fermé pour $\mathring{\ll}$, « mais bien
souvent, $\mathcal{M}'$ ne le sera pas ».

\subsection*{Homomorphismes d'ateliers}

On cherche (pages 62 à 64) des applications $f : \mathfrak{F} \to
\mathfrak{F}'$ entre ateliers qui respectent aussi $\leq$. 1°) Commuter aux Sup
quelconques : $f$ est alors déterminée par sa restriction croissante
$f_{\mathcal{M}} : \mathcal{M} \to \mathfrak{F}'$. 2°) Commuter aux Inf finis,
$f(F \cap G) = f(F) \cap f(G)$, ce qui se ramène à
\[
f(X) \cap f(Y) = \mathrm{Sup}_{Z \in \widetilde{X} \cap \widetilde{Y}} f(Z),
\]
qui « rappelle » la formule $X_{i} \cap X_{j} = \bigcup_{k \leq i,j} X_{k}$ des
pages 45 à 49. Pour donner un sens à la compatibilité avec $\mathring{\ll}$ et
$\mathrel{|\circ|}$, il faut encore :
\begin{quote}
$(**)$ pour toute figure $F$, l'application $G \mapsto f(G)$ est un
isomorphisme de l'ensemble ordonné des sous-figures de $F$ sur celui des
sous-figures de $f(F)$.
\end{quote}
Mais alors $f$ envoie une figure élémentaire sur une figure élémentaire, et l'on
retombe sur $f : \mathcal{M} \to \mathcal{M}'$ comme aux pages 35 et 36 ; les
exemples qui n'envoient pas $\mathcal{M}$ dans $\mathcal{M}'$, comme le produit
de la page 58, lui paraissent « hétéroclites » et sont laissés de côté.

Un tel $f$ n'est pas nécessairement injectif (page 65) : pour un magasin
quasi-ensembliste non ensembliste, l'homomorphisme canonique vers les figures
de ses lieux ne l'est pas. Et même injectif, son image n'est un sous-magasin que
si
\[
(**') \qquad f(X) \ll f(Y) \Longrightarrow X \ll Y ,
\]
ce qui n'est pas automatique ; $\mathring{\ll}$ est alors induite, « qu'en
est-il de $\leq$, $\mathrel{|\circ|}$ ? Il semble qu'il faille le poser en
sus… ».

\pagerange{66}{72}
\section*{Fonctions support exactes (pages 66 à 72)}

Un homomorphisme $f : \mathcal{M} \to \mathrm{Figél}(P)$ donne $|X| = |f(X)|
\subset P$ et en est déterminé, $f(X) = \{|Y| \mid Y \leq X\}$. « Nous voilà
revenus à la situation » de la page 41 :

\begin{quote}
\textbf{Théorème} (page 69). Les homomorphismes $\mathcal{M} \to
\mathrm{Figél}(P)$ correspondent bijectivement aux applications $X \mapsto |X|$
de $\mathcal{M}$ dans $\mathfrak{P}(P)$ telles que :
\end{quote}
\begin{enumerate}
\item[(i)] $X \ll Y \Rightarrow |X| \subset |Y|$ ;
\item[(ii)] $X \mathrel{|\circ|} Y \Rightarrow |X|^{\circ} \cap |Y|^{\circ} =
\emptyset$ ;
\item[(iii)] $|X|^{\circ} \neq \emptyset$ ;
\item[(iv)] $|X| = \bigcup_{Y \leq X} |Y|^{\circ}$.
\end{enumerate}
Une telle application est une \emph{fonction support exacte}.


Que $f(X)$ soit une figure élémentaire indexée par $\widetilde{X}$ est la
proposition des pages 45 à 50, les intérieurs étant disjoints par Mag 4 et (ii).
Le point délicat est $X \mathrel{\mathring{\ll}} Y \Rightarrow |X|^{\circ}
\subset |Y|^{\circ}$ (pages 68 et 69) : $|X| \subset |Y|$ par (i), et pour $Z <
Y$, Mag 4 donne $Z \mathrel{|\circ|} Y$, Mag 3 $X \mathrel{|\circ|} Z$, (ii)
$|X|^{\circ} \cap |Z|^{\circ} = \emptyset$ ; par (iv), $|\partial Y| =
\bigcup_{Z < Y} |Z|^{\circ}$, d'où la conclusion : « on a gagné ! »\footnote{Le
théorème porte deux astérisques que la page n'explique pas ; il renvoie aux
conditions « de la page 59 », sa pagination, notre page 67. La condition (i)
est écrite « $|X| \leq |Y|$ ».}.

La fonction support est \emph{normalisée} si les $|X|$ recouvrent $P$ et en
séparent les points (ou, ce qui revient au même, les $|X|^{\circ}$) ; on s'y
ramène par passage au quotient. Elle est \emph{fidèle} si
\[
(\text{ii fid}) \qquad X \mathrel{|\circ|} Y \iff |X|^{\circ} \cap |Y|^{\circ}
= \emptyset ,
\]
et la marge appelle \emph{réalisation} une fonction support fidèle qui reflète
de plus $\leq$ et $\mathring{\ll}$. Tous les exemples « provenant des contextes
ensemblistes, de topologie ou de topologie modérée » admettent une fonction
support fidèle, de sorte que les propriétés qui en découlent « pourraient être
posées en axiome sans inconvénient pratique ». Seule exception entrevue : les
géométries où la disjonction est définie par une distance $\geq \varepsilon$ ou
$> \varepsilon$.

\emph{Exemples} (pages 71 et 72). 1) $\mathrel{|\circ|}$ est définie par
$\mathring{\ll}$, au sens de la page 39, si et seulement si $X \mapsto
\mathrm{Omb}(X)$, à valeurs dans $\mathfrak{P}(\mathcal{M})$, est une fonction
support fidèle\footnote{On a bien $|\mathrm{Omb}(X)|^{\circ} = \mathrm{Omb}(X)
\smallsetminus \bigcup_{Y < X} \mathrm{Omb}(Y) = \mathrm{Omb}^{\circ}(X)$ par
Mag 2, et (ii fid) devient : $X \mathrel{|\circ|} Y$ si et seulement si $X$ et
$Y$ n'ont pas de $\mathring{\ll}$-minorant commun.}. 2) Les magasins de
« pseudo-cylindres » associés à une corde topologique ou à un ensemble ordonné
ont une fonction support évidente. 3) Pour un magasin trivial, $|X| =
\widetilde{X}$ --- l'adhérence de $X$ pour la topologie de $\mathcal{M}$
associée à $\leq$ --- est une fonction support fidèle, avec $|X|^{\circ} =
\{X\}$.

\pagerange{73}{80}
\section*{Supports intérieurs (pages 73 à 80)}

« Il est temps de revenir aux supports internes dans un magasin. » Pour des
parties $A$, $B$ de $\mathcal{M}$, on écrit $A \mathrel{|\circ|} B$ si tout
élément de $A$ est intérieurement disjoint de tout élément de $B$, et l'on pose
\[
\mathrm{cosupp}^{\circ}(A) = \{ Y \in \mathcal{M} \mid A \mathrel{|\circ|} Y
\}, \qquad \mathrm{supp}^{\circ}(A) =
\mathrm{cosupp}^{\circ}(\mathrm{cosupp}^{\circ}(A)).
\]
Un \emph{support} (intérieur) est une partie $A$ telle que $A =
\mathrm{supp}^{\circ}(A)$ ; $\Sigma_{\mathcal{M}}$ est leur ensemble.

C'est la correspondance de Galois --- la polarité --- associée à la relation
symétrique $\mathrel{|\circ|}$, et $\mathrm{supp}^{\circ}$ est la fermeture
qu'elle définit\footnote{G. Birkhoff, \emph{Lattice Theory} (1940), chapitre
sur les polarités. Une relation symétrique et antiréflexive est ce qu'on
appelle une relation d'orthogonalité, et ses parties « bi-orthogonalement
fermées » forment un orthotreillis complet (espaces d'orthogonalité, J. R.
Dacey 1968, puis D. Foulis et C. Randall). Les rapprochements sont les nôtres ;
la page ne cite personne.}. D'où ce que dit la page 75 :
$\Sigma_{\mathcal{M}}$ est ordonné par inclusion, de plus petit élément
$\emptyset$ et de plus grand $\mathcal{M}$ ; ses Inf sont les intersections ;
$S \mapsto \mathrm{cosupp}^{\circ}(S)$ est une anti-involution, qui échange Inf
et Sup ; et
\[
\mathrm{Sup}_{i}\, S_{i} = \mathrm{supp}^{\circ}\Bigl(\bigcup_{i}
S_{i}\Bigr), \qquad \text{à ne pas confondre avec } \bigcup_{i} S_{i}.
\]
L'antiréflexivité donne de plus $S \cap \mathrm{cosupp}^{\circ}(S) =
\emptyset$, ce qui fait de $\Sigma_{\mathcal{M}}$ un orthotreillis complet ---
qui n'est pas distributif en général : c'est la « pathologie » que les pages 85
à 87 font disparaître. Un support est fermé pour $\mathring{\ll}$ vers le bas
(par Mag 3), mais pas nécessairement pour $\leq$, « et encore moins pour
$\ll$ ».

La page 74 compare avec la notion antérieure de support, $\mathrm{supp}\,X =
\mathrm{supp}^{\circ}(\widetilde{X})$ : le traitement précédent « était vicié
par le fait qu'il se brouillait » avec les supports des parties fermées ; tous
les supports s'expriment par $\mathrm{supp}^{\circ}$, non l'inverse.
\emph{Exemple instructif} : pour un magasin trivial, $\mathrel{|\circ|}$ est la
différence, et toute partie est un support --- à visualiser comme la réunion
des « cellules ouvertes » $X^{\circ}$, $X \in A$. Avec $\parallel$ au contraire,
les supports ne seraient que certaines parties fermées : « jamais par elles on
ne pourra exprimer des supports internes ! »

\subsection*{Supports et lieux (pages 75 à 77)}

Sur les lieux, avec $\parallel$, on a de même $\Sigma_{\mathcal{L}}$. Si le
magasin est strictement divisible, $\mathrm{cosupp}^{\circ}(A) =
\mathrm{cosupp}^{\circ}(\mathrm{Omb}^{\circ}(A))$ pour toute partie $A$, où
$\mathrm{Omb}^{\circ}(A) = \bigcup_{X \in A} \mathrm{omb}(X)^{\circ}$, et
\[
\mathrm{cosupp}^{\circ}_{\mathcal{M}}(A) \cap \mathcal{L} =
\mathrm{cosupp}^{\circ}_{\mathcal{L}}(\mathrm{Omb}^{\circ}(A)),
\]
de sorte que\footnote{La page écrit $\mathrm{Omb}(A)$ dans le second membre de
cette formule ; il faut $\mathrm{Omb}^{\circ}$, comme dans la ligne qui la
précède.} $S \mapsto S \cap \mathcal{L}$ envoie $\Sigma_{\mathcal{M}}$ dans
$\Sigma_{\mathcal{L}}$. En sens inverse, $A \mapsto
\mathrm{supp}^{\circ}_{\mathcal{M}}(A)$, dont la page conjecture qu'il vaut
$\{X \mid \mathrm{omb}(X)^{\circ} \subset A\}$.

\begin{quote}
\textbf{Proposition} (pages 76 et 77). Si $\mathcal{M}$ est strictement
divisible, $S \mapsto S \cap \mathcal{L}$ et $A \mapsto \{X \in \mathcal{M}
\mid \mathrm{omb}(X)^{\circ} \subset A\}$ sont des bijections réciproques entre
$\Sigma_{\mathcal{M}}$ et $\Sigma_{\mathcal{L}}$, qui échangent les
complémentaires $\mathrm{cosupp}^{\circ}$.
\end{quote}

« Il faudra bien que je me tape la vérification ! » La voici\footnote{La
vérification n'est pas sur la page ; elle est nôtre. Strictement divisible
signifie ici que $X \mathrel{|\circ|} Y$ si et seulement si tout lieu intérieur
à $X$ est disjoint de tout lieu intérieur à $Y$ (Mag L 2 avec Mag 3). Pour $A
\in \Sigma_{\mathcal{L}}$, notons $A^{\perp} = \mathrm{cosupp}^{\circ}_{\mathcal{L}}(A)$.
Alors $\mathrm{cosupp}^{\circ}_{\mathcal{M}}(A) = \{Y \mid
\mathrm{omb}(Y)^{\circ} \subset A^{\perp}\}$, et $X$ est dans
$\mathrm{supp}^{\circ}_{\mathcal{M}}(A)$ si et seulement si ses lieux intérieurs
sont disjoints de ceux de tous ces $Y$, parmi lesquels les lieux $y \in
A^{\perp}$ eux-mêmes ($\mathrm{omb}(y)^{\circ} = \{y\}$) : si et seulement si
$\mathrm{omb}(X)^{\circ} \subset A^{\perp\perp} = A$. Cela établit l'égalité
conjecturée et $\mathrm{supp}^{\circ}_{\mathcal{M}}(A) \cap \mathcal{L} = A$.
Inversement, pour $S = \mathrm{cosupp}^{\circ}_{\mathcal{M}}(B)$, on a $S = \{X
\mid \mathrm{omb}(X)^{\circ} \subset \mathrm{Omb}^{\circ}(B)^{\perp}\}$ et $S
\cap \mathcal{L} = \mathrm{Omb}^{\circ}(B)^{\perp}$, d'où $S = \{X \mid
\mathrm{omb}(X)^{\circ} \subset S \cap \mathcal{L}\}$. Enfin
$\mathrm{Omb}^{\circ}(S) = S \cap \mathcal{L}$ pour un support $S$, fermé pour
$\mathring{\ll}$, ce qui donne la compatibilité aux complémentaires.}.

\begin{quote}
\textbf{Corollaire} (page 77). Dans le cas quasi-ensembliste, où $\parallel$ est
la différence sur $\mathcal{L}$, $\Sigma_{\mathcal{M}} \simeq
\mathfrak{P}(\mathcal{L})$, et $\mathrm{supp}^{\circ}(A)$ s'identifie à
$\mathrm{Omb}^{\circ}(A)$.
\end{quote}

« Si je me fatigue à faire une théorie “abstraite”, c'est bien sûr pour
pouvoir parler des supports aussi en dehors des cas quasi-ensemblistes (p. ex.
pour les cordes etc.) » (page 78).

\subsection*{Supports et fonction support (pages 78 à 80)}

Soit une fonction support exacte, normalisée et fidèle $X \mapsto |X|$, et
$\varphi(A) = |A|^{\circ} = \bigcup_{X \in A} |X|^{\circ}$. Par fidélité,
$\mathrm{cosupp}^{\circ}(A) = \{X \mid |X|^{\circ} \cap \varphi(A) =
\emptyset\}$, et tout support $S$ se retrouve à partir de $\varphi(S)$ par $X
\in S \iff |X|^{\circ} \subset \varphi(S)$ (page 79). Donc
\[
\Sigma_{\mathcal{M}} \longrightarrow \mathfrak{P}(P), \qquad S \longmapsto
|S|^{\circ},
\]
est injective, avec pour inverse à gauche $Q \mapsto
\mathrm{supp}^{\circ}(\{X \mid |X|^{\circ} \subset Q\})$. Il voudrait que
l'ordre soit induit et que Sup, Inf et complémentaire soient respectés : « Est-ce
trop demander ? » Oui en général ; d'où la décision de la page 80 de « se
limiter à des supports de figures, finis », les supports « constructibles ».
Le croquis de la marge montre deux parties disjointes dont les supports se
rencontrent.

\pagerange{81}{88}
\section*{Supports constructibles et pleine fidélité (pages 81 à 88)}

Il veut d'abord $|\mathrm{supp}^{\circ} X|^{\circ} = |X|^{\circ}$, c'est-à-dire :
si $|Y|^{\circ} \not\subset |X|^{\circ}$, il existe $Z \mathrel{|\circ|} X$ dont
l'intérieur rencontre celui de $Y$. C'est clair si $|Y|^{\circ}$ rencontre
$|\partial X|$ ; hors de $|X|$, il faut une hypothèse (page 82) : tout point
hors de $|X|$ est dans l'intérieur d'un $Z$ intérieurement disjoint de $X$. Elle
assure (page 83) $|\mathrm{supp}^{\circ} X|^{\circ} = |X|^{\circ}$ et
$|\mathrm{cosupp}^{\circ} X|^{\circ} = P \smallsetminus |X|^{\circ}$, et, « dans
toute sa force », $|\mathrm{supp}\, X|^{\circ} = |X|$\footnote{La fin de la
page 83 écrit « $|X| \cap |Z|^{\circ} \neq \emptyset$, $|Y|^{\circ} \cap
|Z|^{\circ}$ » ; il faut lire $|X| \cap |Z|^{\circ} = \emptyset$ et $|Y|^{\circ}
\cap |Z|^{\circ} \neq \emptyset$.}.

\begin{quote}
\textbf{Définitions} (pages 84 et 85). Soit $\mathfrak{F}$ un atelier. Un
support est \emph{constructible} s'il est de la forme
$\mathrm{supp}^{\circ}(\Phi)$ pour une partie $\Phi$ d'une figure $F \in
\mathfrak{F}$. Une fonction support fidèle est \emph{pleinement fidèle} si pour
toute figure $F \in \mathfrak{F}$,
\[
P \smallsetminus |F| = \bigcup_{Z \mathrel{|\circ|} F} |Z|^{\circ},
\]
c'est-à-dire si tout point hors de $|F|$ est dans l'intérieur d'un $Z$ dont
l'intérieur évite $|F|$\footnote{La page écrit $X \smallsetminus |F|$, et la
marge corrige : « plutôt $P$ ». Sous cette forme la propriété entraîne la
fidélité, dit une autre marge.}.
\end{quote}

Le mot « constructible » est celui de la géométrie algébrique, et l'analogie
est juste : ce sont des réunions finies de strates ouvertes d'une même
stratification\footnote{Les ensembles constructibles de Chevalley, réunions
finies de parties localement fermées, dont EGA fait un usage systématique. Le
rapprochement est le nôtre.}.

\begin{quote}
\textbf{Proposition} (page 85). Si la fonction support est pleinement fidèle
et $\Phi \subset F \in \mathfrak{F}$, alors $|\mathrm{supp}^{\circ}(\Phi)|^{\circ}
= |\Phi|^{\circ}$.
\end{quote}

Il faut voir qu'un $Y \in \mathrm{supp}^{\circ}(\Phi)$ a son intérieur dans
$|\Phi|^{\circ}$. Sinon un point de $|Y|^{\circ}$ est ou bien dans $|F|
\smallsetminus |\Phi|^{\circ} = |F \smallsetminus \Phi|^{\circ}$, donc dans
l'intérieur d'un $Z \in F \smallsetminus \Phi$, intérieurement disjoint de
$\Phi$ (proposition 2 de la page 9) ; ou bien hors de $|F|$, et la pleine
fidélité fournit le $Z$. Dans les deux cas $Z \in
\mathrm{cosupp}^{\circ}(\Phi)$ et $Z$ n'est pas intérieurement disjoint de $Y$.

\begin{quote}
\textbf{Corollaires} (page 86). L'application $\Phi \mapsto
\mathrm{supp}^{\circ}(\Phi)$, des parties de $F$ dans $\Sigma_{\mathcal{M}}$,
commute aux Inf de familles non vides et aux Sup, et
$\mathrm{supp}^{\circ}(\Psi \smallsetminus \Phi) = \mathrm{supp}^{\circ}\Psi \cap
\mathrm{cosupp}^{\circ}(\mathrm{supp}^{\circ}\Phi)$ pour $\Phi \subset \Psi$.
Les supports constructibles subordonnés à $F$ forment une partie de
$\Sigma_{\mathcal{M}}$ stable par ces opérations, sur laquelle $A \mapsto
|A|^{\circ}$ les respecte.
\end{quote}

\begin{quote}
\textbf{Scholie} (pages 86 et 87). Pour les familles \emph{admissibles} de
supports constructibles --- subordonnés à une même figure --- Sup, Inf et
complémentaire relatif ont toutes les propriétés ensemblistes habituelles, et
$A \cap B = \emptyset$ si et seulement si $A \mathrel{|\circ|} B$. « Toute
pathologie disparaît. »
\end{quote}

La page 87 pose alors une question --- un magasin muni d'une fonction support
pleinement fidèle satisfait-il nécessairement le scholie ? ---, à laquelle la
marge répond « Non »\footnote{Telle que la transcription la lit, la question
porte sur un magasin qui \emph{a} une fonction support pleinement fidèle, et la
proposition qui précède y répond par l'affirmative ; la réponse « Non » de la
marge porte donc vraisemblablement sur une autre lecture de la question
(l'existence d'une telle fonction étant à démontrer). Le passage est
incertain.}, et une seconde : existe-t-il des fonctions support fidèles qui ne
soient pas pleinement fidèles ?

\begin{quote}
\textbf{Axiome des supports} (supp 1, page 87). Si $F \in \mathfrak{F}$, $\Phi
\subset F$ et $\Phi' = F \smallsetminus \Phi$, alors
\[
\mathrm{supp}^{\circ}(F) \cap \mathrm{cosupp}^{\circ}(\Phi) =
\mathrm{supp}^{\circ}(\Phi') ,
\]
c'est-à-dire : si $Z \in \mathrm{supp}^{\circ}(F)$ et $Z \mathrel{|\circ|}
\Phi$, alors $Z \in \mathrm{supp}^{\circ}(\Phi')$\footnote{La reformulation de
la page porte « $Z \parallel \Phi$ » ; c'est $\mathrel{|\circ|}$ qu'impose la
première ligne.}.
\end{quote}

« C'est le premier axiome de nature délicate, sur lequel nous tombions. » Il
ne fait intervenir que $\leq$ et $\mathrel{|\circ|}$. Il est satisfait chaque
fois qu'il existe une fonction support pleinement fidèle --- c'est le dernier
corollaire. La page ajoute qu'il l'est dans les magasins fidèles « car Omb y est
une fonction support extérieure pleinement fidèle » ; c'est trop
dire\footnote{Dans le magasin des intervalles d'une droite (pages 89 à 93),
fidèle d'après la page 93, prenons $F = \widetilde{S}_{0,2}$ et $W = S_{1,3}$.
$W$ n'est pas dans $\mathrm{Omb}(F)$, mais il n'est dans
$\mathrm{Omb}^{\circ}(Z)$ que pour des $Z$ qui contiennent $[1,3]$, et qui
chevauchent donc $S_{0,2}$ : $\mathrm{Omb}$ n'est pas pleinement fidèle.
L'axiome y est pourtant vrai, parce que la fonction support ensembliste $X
\mapsto |X| \subset \mathbb{R}$, elle, est pleinement fidèle. Que tout magasin
fidèle satisfasse l'axiome reste donc à voir.}.

Programme de la page 88 : 1°) les magasins définis par un ensemble ordonné ;
2°) les magasins définis par $\leq$ et $\mathrel{|\circ|}$ seuls,
$\mathring{\ll}$ s'en déduisant ; 3°) voir si dans ces magasins
« disjonctifs » l'axiome des supports peut être en défaut --- « Oui ! », dit la
marge. Seul 1°) est fait dans ce dossier.

\pagerange{89}{96}
\section*{Le magasin d'un ensemble ordonné (pages 89 à 96)}

Le premier juillet. Soit $L$ un ensemble ordonné. On pose
\[
\mathcal{M} = \mathcal{M}_{0} \amalg \mathcal{M}_{1}, \qquad
\mathcal{M}_{0} = \{P_{x} \mid x \in L\}, \qquad \mathcal{M}_{1} = \{S_{a,b}
\mid a, b \in L,\ a < b\},
\]
les \emph{points} et les \emph{intervalles}\footnote{La page note $P_{x} =
\{\{x\}\}$ et $S_{a,b} = \{[a,b], \{a\}, \{b\}\}$ comme figures ensemblistes ;
la marge observe qu'il serait plus simple de prendre $\mathcal{M}_{0} = L$ et
$\mathcal{M}_{1}$ l'ensemble des couples $a < b$, et la page 103 finit par
identifier $P_{a}$ à $a$. À la page 119, il préfère appeler $S_{a,b}$ un
« intervalle » et réserver « segment » à $\widetilde{S}_{a,b}$ ; nous le
faisons dès maintenant.}, avec
\[
\begin{aligned}
X \leq Y &\iff X = Y, \text{ ou } X = P_{x},\ Y = S_{a,b},\ x \in \{a, b\}
\qquad (\partial S_{a,b} = \{P_{a}, P_{b}\}) ; \\
X \mathrel{\mathring{\ll}} Y &\iff
\begin{cases}
P_{x} = P_{y}, \\
X = P_{x},\ Y = S_{a,b},\ a < x < b, \\
X = S_{a,b},\ Y = S_{a',b'},\ a' \leq a < b \leq b' ;
\end{cases} \\
X \mathrel{|\circ|} Y &\iff
\begin{cases}
X = P_{x},\ Y = P_{y} : & x < y \text{ ou } y < x, \\
X = P_{x},\ Y = S_{a,b} : & x \leq a \text{ ou } b \leq x, \\
X = S_{a,b},\ Y = S_{a',b'} : & b \leq a' \text{ ou } b' \leq a .
\end{cases}
\end{aligned}
\]
Si $L$ est totalement ordonné et \emph{divisible} --- entre deux éléments il y
en a un troisième ---, ce sont les relations ensemblistes entre les parties
$\{x\}$ et $[a,b]$ de $L$, avec $|S_{a,b}|^{\circ} = \,]a,b[$ ; c'est le
\emph{magasin unidimensionnel} de $L$\footnote{Dans la deuxième ligne de
$\mathring{\ll}$ la page porte « $x \in |X|^{\circ}$ » : lire $|Y|^{\circ}$.}.
Dans le cas général (page 90), on prend les mêmes formules pour définitions ;
$|S_{a,b}|^{\circ} = \,]a,b[$ peut être vide, et $\mathrel{|\circ|}$ entraîne
la disjonction des intérieurs ensemblistes sans lui être équivalente\footnote{La
page 91 écrit, pour le cas d'un point et d'un intervalle, « $x \leq a$ ou $y
\geq b$ » : lire $x \geq b$.}. On vérifie Mag 1 à Mag 4 (page 91), et
\[
X \ll Y \iff |X| \subset |Y| ;
\]
les lieux sont les $P_{x}$, et l'ensemble des lieux s'identifie à $L$.

\begin{quote}
\textbf{Propriétés} (pages 92 et 93). a) Mag L 1 équivaut à la divisibilité de
$L$. b) Mag L 2 équivaut à : si $x$ est comparable et distinct de tout point de
$]a,b[$, alors $x \leq a$ ou $x \geq b$. c) $\mathcal{M}$ est fidèle si et
seulement si $L$ est totalement ordonné.
\end{quote}

Pour b), la page donne une forme équivalente en termes des
$L_{\leq c}$\footnote{« (iii) a) pour $a$ non maximal, $\bigcap_{c > a}
L_{\leq c} = L_{\leq a}$ ; b) l'énoncé dual ». Les conditions (i) et (ii) qui
précèdent sont en partie illisibles et l'équivalence n'est pas démontrée ; nous
ne la garantissons pas. La condition de l'énoncé b) est, elle, toujours vraie si
$L$ est totalement ordonné.}, et note qu'elle n'est « le plus souvent pas
satisfaite » pour les ensembles ordonnés tirés d'espaces topologiques (croquis
d'un fuseau traversé par deux bandes). Pour c) (page 93) : il faut que deux
points distincts soient comparables, et cela suffit, car si deux intervalles
$S_{a,b}$, $S_{a',b'}$ ne sont pas intérieurement disjoints,
$S_{\max(a,a'),\min(b,b')}$ est intérieur aux deux. Le corollaire qui suit dit
qu'un $L$ totalement ordonné et divisible donne un magasin
ensembliste\footnote{La page qualifie un tel magasin d'un mot lu
« parfait », avec doute.}.

\subsection*{Compatibilité, drapeaux, figures canoniques (pages 94 à 96)}

Un point est compatible avec $Y$ si et seulement s'il lui est égal ou
intérieurement disjoint (« vrai pour des lieux dans un magasin ») ; deux
intervalles sont compatibles si et seulement s'ils sont égaux ou intérieurement
disjoints, ce qui comprend les intervalles contigus $S_{a,b}$, $S_{b,c}$. Soit
$\mathrm{Drap}(L)$ l'ensemble des \emph{drapeaux} de $L$, ses parties finies
totalement ordonnées\footnote{C'est l'ensemble des simplexes du complexe des
chaînes de $L$ (son « complexe d'ordre ») ; le nom est le nôtre.}, et
$\mathrm{Drap}^{*}(L)$ celui des drapeaux non vides. Avec $\delta P_{x} = \{x\}$,
$\delta S_{a,b} = \{a, b\}$, et pour un drapeau $T = \{t_{1} < \cdots < t_{n}\}$
ses \emph{intervalles interstitiels} $S_{t_{i}, t_{i+1}}$ et ses \emph{lieux
interstitiels} $P_{t_{i}}$ :

\begin{quote}
\textbf{Proposition} (page 96). Une partie finie non vide $F$ de $\mathcal{M}$
est formée d'éléments deux à deux compatibles si et seulement si $\delta F =
\bigcup_{X \in F} \delta X$ est un drapeau et si les éléments de $F$ sont
interstitiels pour $\delta F$.
\end{quote}

La démonstration est indiquée par récurrence sur le cardinal. Pour un drapeau
$T$, la plus grande telle partie est la \emph{figure canonique} $F_{T}$, formée
de tous les points et intervalles interstitiels de $T$ ; elle détermine $T =
\delta F_{T}$. Un début de calcul des supports de $\widetilde{S}_{a,b}$ est
biffé au bas de la page 96.

\pagerange{97}{102}
\section*{Figures et supports élémentaires (pages 97 à 102)}

Le 2 juillet. Sur $\mathrm{Drap}^{*}(L)$, on pose $T < T'$ si $t < t'$ pour
tous $t \in T$, $t' \in T'$ ; on peut alors parler de drapeaux de drapeaux.

\begin{quote}
\textbf{Lemme et théorème} (pages 97 et 98). $F_{T} \parallel F_{T'}$ si et
seulement si $T < T'$ ou $T' < T$. L'application
\[
\mathcal{T} \longmapsto F_{\mathcal{T}} = \bigcup_{T \in \mathcal{T}} F_{T},
\qquad \mathrm{Drap}(\mathrm{Drap}^{*}(L)) \longrightarrow
\{\text{figures finies de } \mathcal{M}\},
\]
est bijective.
\end{quote}

Une figure finie se décompose en « composantes connexes », au sens combinatoire
évident, qui sont des $F_{T}$ et « se suivent à la queue-leu-leu » (croquis de
la page 99). Une marge en donne une autre présentation : une figure est un
drapeau $\mathfrak{D} = \delta F$ et une partie de l'ensemble de ses $n - 1$
intervalles interstitiels, soit $2^{n-1}$ figures pour un drapeau à $n$
éléments.

Pour les supports, il faut une hypothèse (page 99) :
\begin{quote}
\textbf{Hyp. H.} a) Ou bien $L$ a un plus petit élément, ou bien tout élément
de $L$ en majore strictement un autre. b) L'énoncé dual.
\end{quote}

\begin{quote}
\textbf{Lemme 1} (page 99). Sous H, pour tout $X \in \mathcal{M}$,
$\mathrm{supp}(X) = \mathrm{supp}^{\circ}(\widetilde{X}) = \mathrm{Omb}(X) =
\{Z \mid |Z| \subset |X|\}$. Explicitement :
\[
\mathrm{supp}(P_{a}) = \{P_{a}\}, \qquad \mathrm{supp}(S_{a,b}) = \{P_{x} \mid
a \leq x \leq b\} \cup \{S_{x,y} \mid a \leq x < y \leq b\},
\]
\[
\mathrm{supp}^{\circ}(S_{a,b}) = \mathrm{Omb}^{\circ}(S_{a,b}) =
\mathrm{supp}(S_{a,b}) \smallsetminus \{P_{a}, P_{b}\},
\]
\[
\mathrm{supp}^{\circ}(\{S_{a,b}, P_{a}\}) = \mathrm{Omb}^{\circ}(S_{a,b}) \cup
\{P_{a}\}.
\]
\end{quote}

L'inclusion $\mathrm{Omb}(X) \subset \mathrm{supp}(X)$ est triviale. Pour
l'autre, cas de $P_{a}$ (pages 100 et 101) : si $a$ est un plus petit ou un
plus grand élément, tout autre élément de $\mathcal{M}$ est intérieurement
disjoint de $P_{a}$, et $\mathrm{supp}^{\circ}(P_{a}) = \{P_{a}\}$ par
antiréflexivité ; sinon, H fournit $x < a < y$, et $X \in
\mathrm{supp}^{\circ}(P_{a})$ doit être intérieurement disjoint de $S_{x,a}$ et
de $S_{a,y}$, ce qui force $\delta X \geq a$ et $\delta X \leq a$, donc $X =
P_{a}$\footnote{L'argument exclut $\delta X \leq x$ parce qu'on aurait alors $X
\in \mathrm{cosupp}^{\circ}(P_{a})$, et $\mathrm{cosupp}^{\circ}(\Phi) \cap
\mathrm{supp}^{\circ}(\Phi) = \emptyset$ pour toute partie $\Phi$. Il n'utilise
pas que $L$ soit totalement ordonné.}. Le cas de $S_{a,b}$ est analogue.

\emph{Le « canular »} (page 102). Pour voir ce qui se passe sans H, la page
prend $L = \{a, b, c\}$ avec $a < b$ et $c < b$ seulement, et calcule
$\mathrm{cosupp}^{\circ}(P_{a}) = \{P_{b}, S_{a,b}, S_{c,b}\}$ et
$\mathrm{supp}^{\circ}(P_{a}) = \{P_{a}, P_{b}\}$ ; la marge écrit
« Canular ». Le calcul est faux\footnote{$P_{a}$ n'est pas intérieurement
disjoint de $S_{c,b}$ ($a$ n'est ni $\leq c$ ni $\geq b$), et $P_{b}$, qui n'est
pas disjoint de lui-même, ne peut être dans $\mathrm{supp}^{\circ}(P_{a})$ ; en
fait $\mathrm{cosupp}^{\circ}(P_{a}) = \{P_{b}, S_{a,b}\}$ et
$\mathrm{supp}^{\circ}(P_{a}) = \{P_{a}\}$.}, mais la suite est juste :
$\mathrm{cosupp}^{\circ}(\widetilde{S}_{a,b}) = \emptyset$, donc
$\mathrm{supp}^{\circ}(\widetilde{S}_{a,b}) = \mathcal{M} \neq
\mathrm{Omb}(S_{a,b})$, l'ensemble à cinq éléments contenant en plus $S_{c,b}$ et
$P_{c}$. Le lemme 1 tombe donc en défaut pour ce $L$, qui ne satisfait pas H
($a$ est minimal sans être le plus petit). Un second exemple, en losange, est
esquissé et laissé là, et un lemme 2 sur les supports d'une réunion est
biffé.

\pagerange{102}{110}
\section*{Préfigures (pages 102 à 110)}

Le 4 juillet : « Détermination des profigures, ou ensembles constructibles dans
$\mathcal{M}$ »\footnote{Les pages 102 à 112 disent le plus souvent
« profigure », les pages 88, 103 (dans « Préfigcomm ») et 117 à 126
« préfigure » ; c'est la même notion, et le chapitre IX garde « préfigure ».}.
Une \emph{préfigure} est une partie d'une figure ; sa \emph{figure engendrée}
$\overline{\Phi}$ est son adhérence pour $\leq$, et $\delta\Phi = \delta
\overline{\Phi} = \overline{\Phi} \cap \mathcal{M}_{0}$ est l'ensemble de ses
\emph{sommets}, un drapeau, vide si et seulement si $\Phi = \emptyset$.

\emph{Préfigures connexes} (page 103). Ce sont : a) les figures connexes
$F_{T}$, $T \in \mathrm{Drap}^{*}(L)$ ; b) et c) les $F_{T} \smallsetminus
\alpha$, où $\alpha$ est une partie du bord $\partial T$ de $T$ --- son plus
petit et son plus grand élément si $\mathrm{card}\, T \geq 2$, vide sinon.
Pour $\mathrm{card}\, T = 2$ : l'intervalle fermé, semi-ouvert, ouvert.

\emph{Le bord par composantes} (pages 104 et 105). Pour une préfigure
quelconque, de composantes connexes $\Phi_{i}$, la page pose $\partial_{c}\Phi
= \bigcup_{i} \partial\Phi_{i}$ et $\Phi^{\bullet} = \bigcup_{i}
\Phi_{i}^{\bullet}$, où $\Phi_{i}^{\bullet}$ est $\partial\Phi_{i}$, ou $\{a\}$
si $\Phi_{i} = \{P_{a}\}$\footnote{La page note ce bord $\partial\Phi$ ; nous
écrivons $\partial_{c}\Phi$ pour le distinguer de celui des pages 118 à 126.
La page 105 le caractérise : $a \in \partial_{c}\Phi$ si un élément de $\Phi$
est incident à $a$, et un seul quand $P_{a} \in \Phi$.}. Exemples, pour $a < b <
c < d$ :
\[
\begin{aligned}
\Phi &= \{S_{a,b},\ S_{b,c},\ P_{d}\} : & \delta\Phi &= \{a,b,c,d\}, &
\partial_{c}\Phi &= \{a,b,c\}, & \Phi^{\bullet} &= \{a,b,c,d\} ; \\
\Phi &= \{S_{a,b},\ P_{b},\ S_{b,c},\ P_{d}\} : & \delta\Phi &= \{a,b,c,d\}, &
\partial_{c}\Phi &= \{a,c\}, & \Phi^{\bullet} &= \{a,c,d\}.
\end{aligned}
\]
Dans le premier, $S_{a,b}$ et $S_{b,c}$ sont deux composantes : sans le point
$P_{b}$, elles ne se touchent pas. On a toujours $\partial_{c}\Phi \subset
\Phi^{\bullet} \subset \delta\Phi$, et $\Phi^{\bullet}$ est la réunion disjointe
de $\partial_{c}\Phi$ et des sommets isolés. La marge de la page 104 introduit
déjà l'\emph{ordre} d'un sommet, le nombre d'intervalles de $\Phi$ qui y
aboutissent, qui servira page 118.

\subsection*{L'ordre de position (pages 105 à 108)}

\begin{quote}
\textbf{Définition} (page 105). Sur $\mathcal{M}$, $X <_{\mathrm{pos}} Y$ si
$X$ est entièrement avant $Y$ : $P_{x} <_{\mathrm{pos}} P_{y}$ si $x < y$ ;
$P_{x} <_{\mathrm{pos}} S_{a,b}$ si $x \leq a$ ; $S_{a,b} <_{\mathrm{pos}}
P_{y}$ si $b \leq y$ ; $S_{a,b} <_{\mathrm{pos}} S_{a',b'}$ si $b \leq a'$.
\end{quote}

C'est un ordre strict, et (page 106)
\[
X \mathrel{|\circ|} Y \iff X <_{\mathrm{pos}} Y \text{ ou } Y <_{\mathrm{pos}}
X :
\]
\emph{la disjonction intérieure est la comparabilité pour l'ordre de position}.
Quand $L$ est totalement ordonné et divisible, c'est l'ordre d'intervalles,
au sens d'aujourd'hui, de la famille des intérieurs $\{x\}$ et
$]a,b[$\footnote{Un ordre d'intervalles (P. Fishburn, 1970) est l'ordre sur
une famille d'intervalles d'une droite où $I < J$ si $I$ est entièrement à
gauche de $J$. Le rapprochement est le nôtre.}. Entre préfigures, $\Phi
<_{\mathrm{pos}} \Psi$ si chaque élément de $\Phi$ est avant chaque élément de
$\Psi$ ; et $\Phi <^{1} \Psi$ si de plus, lorsque le dernier sommet de $\Phi$
est le premier de $\Psi$, disons $a$, $P_{a}$ n'est ni dans $\Phi$ ni dans
$\Psi$ --- on exclut le « raccord » où $\Phi$ et $\Psi$ ne seraient pas
ouvertes et fermées dans $\Phi \cup \Psi$.

\begin{quote}
\textbf{Lemme} (pages 107 et 108). Soient $\Phi, \Psi$ des préfigures connexes.
$\Phi \cap \Psi = \emptyset$ et $\Phi \cup \Psi$ est une préfigure si et
seulement si $\Phi <_{\mathrm{pos}} \Psi$ ou $\Psi <_{\mathrm{pos}}
\Phi$\footnote{La page écrit « $\Phi \leq_{\mathrm{pos}} \Psi$ » ; deux
préfigures disjointes ne peuvent être égales, et l'on lit l'alternative
stricte.} ; et $\Phi$, $\Psi$ sont les composantes connexes de $\Phi \cup \Psi$
si et seulement si $\Phi <^{1} \Psi$ ou $\Psi <^{1} \Phi$.
\end{quote}

\begin{quote}
\textbf{Théorème} (pages 108 à 110). 1) $(D, \alpha) \mapsto F_{D}
\smallsetminus \alpha$, pour $D \in \mathrm{Drap}^{*}(L)$ et $\alpha \subset
\partial D$, est une bijection sur les préfigures connexes. 2) Une telle
préfigure est fermée si et seulement si $\alpha = \emptyset$ ; elle est ouverte
si et seulement si $\mathrm{card}\, D \geq 2$ et $\alpha = \partial D
\smallsetminus \Theta$, où $\Theta$ est l'ensemble des éléments extrémaux de
$L$ (sous H). 3) Toute préfigure est, de façon unique, la réunion disjointe
d'une chaîne $\Phi_{1} <^{1} \cdots <^{1} \Phi_{n}$ de préfigures connexes, ses
composantes. 4) Elle est fermée, ou ouverte, si et seulement si ses composantes
le sont.
\end{quote}

« Dans cette description des profigures et des figures, les relations $<^{1}$
n'ont pas intervenu encore, ni (explicitement) la relation $\mathrel{|\circ|}$ »
--- c'est l'ordre de position qui a servi. « Mais maintenant que je passe aux
supports, il n'en sera pas de même. »

\pagerange{110}{117}
\section*{Supports des préfigures (pages 110 à 117)}

La détermination de $\mathrm{supp}^{\circ}\Phi$ se fait en étapes.

\begin{quote}
\textbf{Énoncé réductif} (page 110). Si $\Phi$ a pour composantes connexes les
$\Phi_{i}$, $\mathrm{supp}^{\circ}\Phi = \coprod_{i}
\mathrm{supp}^{\circ}(\Phi_{i})$.

\textbf{Cas connexe} (pages 111 et 112). a) Pour $X \in \mathcal{M}$,
$\mathrm{supp}^{\circ}X = \mathrm{Omb}^{\circ}X$, et pour $\beta \subset
\partial S_{a,b}$, $\mathrm{supp}^{\circ}(\{S_{a,b}\} \cup \beta) =
\mathrm{Omb}^{\circ}(S_{a,b}) \cup \beta$ : c'est le lemme 1 de la page 99. b)
Pour $T \in \mathrm{Drap}^{*}(L)$ de premier élément $a$ et de dernier $b$, et
$\Phi = F_{T} \smallsetminus \alpha$ avec $\alpha \subset \{P_{a}, P_{b}\}$ de
complémentaire $\beta$,
\[
\mathrm{supp}^{\circ}\Phi = \mathrm{Omb}(S_{a,b}) \smallsetminus \alpha =
\mathrm{Omb}^{\circ}(S_{a,b}) \amalg \beta .
\]
\end{quote}
La page écrit, en b), $\mathrm{supp}\, \Phi$ sans l'anneau\footnote{Pour une
préfigure non fermée, c'est $\mathrm{supp}^{\circ}$ qui a un sens ; le premier
membre de la page, $\mathrm{Omb}(\widetilde{S}_{a,b}) \smallsetminus \alpha$,
le confirme.}.

Ce qui est le support d'une préfigure connexe ne dépend que de ses extrémités
et de celles qu'elle contient : les sommets intérieurs ne comptent pas.

L'énoncé réductif résulte, par récurrence sur le nombre de composantes, d'un
lemme (page 112) : \emph{si $\Phi <_{\mathrm{pos}} X <_{\mathrm{pos}} \Psi$
pour un $X$, et $\Phi$, $\Psi$ des préfigures, alors $\mathrm{supp}^{\circ}(\Phi
\cup \Psi) = \mathrm{supp}^{\circ}\Phi \cup \mathrm{supp}^{\circ}\Psi$}. Un tel
$X$ existe exactement quand $\Phi <^{1} \Psi$ : c'est $S_{a,b}$ si le dernier
sommet $a$ de $\Phi$ précède le premier $b$ de $\Psi$, $P_{a}$ si $a = b$. Comme
$X \in \mathrm{cosupp}^{\circ}(\Phi \cup \Psi)$, tout $Z$ du support est avant
ou après $X$, et l'on montre (pages 113 et 114) que dans le premier cas il est
dans $\mathrm{supp}^{\circ}\Phi$, dans le second dans
$\mathrm{supp}^{\circ}\Psi$. « Je n'ai pas eu à utiliser Hyp., mais j'ai utilisé
que $\Phi$, $\Psi$ sont des profigures. »

Le cas b) résulte de a) et de l'énoncé réductif par récurrence sur
$\mathrm{card}\, T$, grâce au
\begin{quote}
\textbf{Lemme de composition des intervalles} (page 114). Pour $a < b < c$,
$\mathrm{supp}^{\circ}(\{S_{a,b}, P_{b}, S_{b,c}\}) =
\mathrm{supp}^{\circ}(S_{a,c})$.
\end{quote}
Les deux parties ont le même $\mathrm{cosupp}^{\circ}$, car être avant ou après
l'une, c'est être avant ou après l'autre --- « trivial (sans Hyp) ».

\begin{quote}
\textbf{Théorème} (pages 115 à 117). 1) Tout support constructible $S$ s'écrit
$\mathrm{supp}^{\circ}\Phi$ avec $\Phi$ une préfigure \emph{non morcelée},
c'est-à-dire sans sommet « redondant » (un point de $\Phi$ où aboutissent deux
intervalles de $\Phi$), et alors
\[
\mathrm{supp}^{\circ}\Phi = \mathrm{Omb}^{\circ}(\Phi) = \coprod_{X \in \Phi}
\mathrm{Omb}^{\circ}(X).
\]
$\Phi$ est déterminée par $S$ : ses intervalles sont les intervalles maximaux de
$S$ pour $\mathring{\ll}$, et ses points sont les extrémités de ces intervalles
qui sont dans $S$, et les points de $S$ maximaux pour $\ll$. 2) Il y a un plus
petit support constructible fermé contenant $S$, $\overline{S} =
\mathrm{supp}^{\circ}(\overline{\Phi})$.
\end{quote}

Une note de la page 116 propose « irredondante » plutôt que « non morcelée ».
Le corollaire qui précède, page 115, est biffé. Pour 2), la page n'est « pas
sûr[e] » de l'expression et note que $\overline{\Phi}$ peut être morcelée quand
$\Phi$ ne l'est pas ; la marge dit « Vérifier ! », ce que nous ne faisons
pas\footnote{Le croquis de la marge : deux intervalles fermés raccordés en un
point, « $\Phi$ non morcelée » quand le point manque, et son adhérence
« morcelée ».}. Le 3), une « description combinatoire des supports
constructibles via les profigures non morcelées », est posé comme question et
fait l'objet de la suite.

\pagerange{117}{126}
\section*{Sommets, et la description des préfigures (pages 117 à 126)}

Toujours le 4 juillet, il veut une présentation « autoduale » des supports
constructibles par les préfigures irredondantes. Il faut pour cela « omettre »
les points isolés et les « points frontière intérieurs » --- les sommets de
$\overline{\Phi}$ qui ne sont pas dans $\Phi$ et ne sont pas des bords.

\subsection*{Zoologie des sommets (pages 118 et 119)}

Pour $a \in L$, l'\emph{ordre} $\mathrm{ord}(a, \Phi) \in \{0, 1, 2\}$ est le
nombre d'intervalles de $\Phi$ incidents à $a$. Un sommet d'ordre 0 est dans
$\Phi$ (lemme de la page 118) : c'est un \emph{sommet isolé}. Les sommets
d'ordre 1 forment le \emph{bord}
\[
\partial\Phi = \{a \in \delta\Phi \mid \mathrm{ord}(a,\Phi) = 1\},
\]
\emph{propre} s'ils sont dans $\Phi$, \emph{impropre} sinon. Les sommets d'ordre
2 sont \emph{internes} : \emph{redondants} (ou « de morcellement ») s'ils sont
dans $\Phi$, \emph{lacunaires} sinon. Dans
\[
\Phi = \{P_{a},\ S_{a,b},\ S_{b,c},\ P_{c},\ S_{c,d},\ P_{e}\}, \qquad a < b < c
< d < e,
\]
il y a un sommet de chacune des cinq espèces : $e$ isolé, $a$ bord propre, $d$
bord impropre, $c$ redondant, $b$ lacunaire, et $\partial\Phi = \{a, d\}$.

Ce bord n'est pas le $\partial_{c}$ des pages 104 et 105 : un sommet lacunaire,
d'ordre 2, n'est pas dans $\partial\Phi$, mais il est au bord de chacune des
deux composantes qu'il sépare, donc dans $\partial_{c}\Phi$. Dans l'exemple
ci-dessus, $\partial_{c}\Phi = \{a, b, d\}$ ; dans le premier exemple de la page
104, $\partial\Phi = \{a, c\}$ et $\partial_{c}\Phi = \{a, b, c\}$. Les deux
coïncident pour une préfigure connexe.

\subsection*{Le complété $\Psi$, et la proposition (pages 120 à 124)}

À une préfigure $\Phi$ on associe (page 120) son adhérence $\overline{\Phi}$,
on en ôte les sommets isolés --- c'est « l'adhérence purement
1-dimensionnelle » ---, puis on efface les sommets redondants de ce qui reste,
en fusionnant les intervalles qui s'y raccordent. On obtient une figure
$\Psi = \Phi_{\mathrm{comp}}$, la \emph{complétion} de $\Phi$, réunion
disjointe de segments $\widetilde{S}_{a_{i}, b_{i}}$ avec
\[
a_{1} < b_{1} < a_{2} < b_{2} < \cdots < a_{n} < b_{n}, \qquad \delta\Psi =
\partial\Psi = \partial\Phi = \{a_{1}, b_{1}, \ldots, a_{n}, b_{n}\}.
\]

\begin{quote}
\textbf{Proposition} (pages 121 à 124). Une préfigure $\Phi$ est déterminée, de
façon unique, par les données suivantes, et toute donnée de ce type en provient :
\end{quote}
\begin{enumerate}
\item[1°)] la figure $\Psi$, ou la suite $a_{1} < b_{1} < \cdots < a_{n} <
b_{n}$ ;
\item[2°)] une partie $\partial_{0} \subset \partial\Psi$, les sommets de bord
propres ;
\item[3°)] un drapeau $\Delta \in \mathrm{Drap}^{*}(L)$ contenant
$\partial\Psi$, l'ensemble des sommets de $\Phi$ ; il se décompose en
\[
\Delta = \partial\Psi \amalg \Delta_{\mathrm{int}} \amalg \Delta_{\mathrm{is}},
\qquad \Delta_{\mathrm{int}} = \Delta \cap \mathrm{Omb}^{\circ}(\Psi),
\]
les sommets intérieurs aux intervalles $S_{a_{i},b_{i}}$ et les sommets
isolés ;
\item[4°)] une partie $\Delta_{\mathrm{lac}} \subset \Delta_{\mathrm{int}}$, les
sommets lacunaires, de complémentaire $\Delta_{\mathrm{red}}$.
\end{enumerate}
On a alors $\Phi_{0} = \partial_{0} \amalg \Delta_{\mathrm{red}} \amalg
\Delta_{\mathrm{is}}$, et $\Phi_{1}$ est l'ensemble des intervalles
interstitiels de $\Delta$ contenus dans l'un des $S_{a_{i}, b_{i}}$ ; et
$\delta\Phi = \Delta$, $\partial\Phi = \partial\Psi$, les sommets isolés,
redondants, lacunaires de $\Phi$ étant $\Delta_{\mathrm{is}}$,
$\Delta_{\mathrm{red}}$, $\Delta_{\mathrm{lac}}$.


\begin{quote}
\textbf{Corollaire 1} (page 124). $\Phi$ est irredondante si et seulement si
$\Delta_{\mathrm{lac}} = \Delta_{\mathrm{int}}$ ; se donner une préfigure
irredondante, c'est se donner $\Psi$, et $\partial_{0} \subset \partial\Psi
\subset \Delta$ avec $\Delta \in \mathrm{Drap}^{*}(L)$.

\textbf{Corollaire 2} (page 125, « pour mémoire »). Sous H, pour $\Phi$
irredondante, $\mathrm{supp}^{\circ}\Phi = \mathrm{Omb}^{\circ}(\Phi_{1}) \amalg
\Phi_{0}$.
\end{quote}

\subsection*{Une présentation plus simple, et le complémentaire (pages 125 et 126)}

« Ça a l'air franchement compliqué comme présentation. Ne vaut-il pas mieux
dire ceci. » Se donner une préfigure, c'est se donner
\begin{enumerate}
\item[1°)] son drapeau de sommets $\Delta = \delta\Phi \in
\mathrm{Drap}^{*}(L)$ ;
\item[2°)] une partie $\Phi_{1}$ de l'ensemble des intervalles interstitiels de
$\Delta$ ;
\item[3°)] une partie $\Phi_{0}$ de $\Delta$ ;
\end{enumerate}
avec la condition que tout sommet de $\Delta$ hors de $\Phi_{0}$ soit
l'extrémité d'un intervalle de $\Phi_{1}$, $\Delta = \Phi_{0} \cup
\bigcup_{X \in \Phi_{1}} \partial X$ --- faute de quoi $\delta(\Phi_{0} \cup
\Phi_{1})$ serait plus petit que $\Delta$ (« exemple : $\Phi_{0} = \Phi_{1} =
\emptyset$ ! »)\footnote{Le nom que la page donne à cette condition, lu
« Pur », est incertain.}.

La \emph{préfigure complémentaire} $\Phi'^{*}$ a les mêmes sommets $\Delta$, et
pour parties $\Phi'^{*}_{0} = \Delta \smallsetminus \Phi_{0}$ et
$\Phi'^{*}_{1}$, les intervalles interstitiels de $\Delta$ qui ne sont pas dans
$\Phi_{1}$. Soient $a$ et $b$ le premier et le dernier élément de $\Delta$, et
$\mathcal{M}_{<a}$, $\mathcal{M}_{>b}$ les éléments de $\mathcal{M}$ avant
$P_{a}$ et après $P_{b}$ pour l'ordre de position ; sous H, $\mathcal{M}_{<a}$
est vide si et seulement si $a$ est le plus petit élément de $L$. Alors
\[
\mathrm{cosupp}^{\circ}(\Phi) = \mathrm{Omb}^{\circ}(\Phi'^{*}) \amalg
\mathcal{M}_{<a} \amalg \mathcal{M}_{>b}, \qquad
\mathrm{Omb}^{\circ}(\Phi'^{*}) = \coprod_{X \in \Phi'^{*}}
\mathrm{Omb}^{\circ}(X) .
\]
Le complémentaire d'un support constructible est, au voisinage de ses sommets,
le support de la préfigure complémentaire. La page écrit
$\mathrm{Omb}^{\circ}(\Phi'^{*}_{1})$ dans le premier membre, l'indice 1
récrit sur un autre chiffre ; il faut la préfigure complémentaire
entière\footnote{Avec $\Phi'^{*}_{1}$ seul, la formule oublierait les points de
$\Delta \smallsetminus \Phi_{0}$. Sur la droite réelle, pour $\Phi =
\{S_{0,1}\}$ (l'intervalle ouvert), $P_{0}$ et $P_{1}$ sont intérieurement
disjoints de $S_{0,1}$, donc dans $\mathrm{cosupp}^{\circ}(\Phi)$, et ne sont
ni dans $\mathcal{M}_{<0}$, ni dans $\mathcal{M}_{>1}$, ni dans
$\mathrm{Omb}^{\circ}(\Phi'^{*}_{1}) = \emptyset$ ; ils sont dans
$\mathrm{Omb}^{\circ}(\Phi'^{*})$, puisque $\Phi'^{*}_{0} = \{P_{0}, P_{1}\}$.
La ligne suivante de la page définit d'ailleurs
$\mathrm{Omb}^{\circ}(\Phi'^{*})$ pour la préfigure entière. On vérifie sans
peine la formule ainsi corrigée quand $L$ est totalement ordonné et
divisible.}.

Le dossier s'arrête là, sur la comparaison de $\mathrm{cosupp}^{\circ}(\Phi)$
et de $\mathrm{Omb}^{\circ}(\Phi'^{*})$. Le chapitre IX (dossier 156-9) la
reprend le 5 juillet, avec les intervalles infinis $I_{<a}$, $I_{>a}$ qui
absorbent $\mathcal{M}_{<a}$ et $\mathcal{M}_{>b}$, sous la forme
$\mathrm{cosupp}^{\circ}(F) = \mathrm{Omb}^{\circ}(F^{*})$.

\end{document}
